% OUTLINE B v6 COMBINED: Full TMLR paper. Built from v5.5 base + all
% recent findings (V8-V10 writeups, factor overlap, per-edge fingerprints,
% attribution-causal orthogonality, background/modulation decomposition,
% per-counterfactual stability, conceptor steering, causal discovery,
% distributed logit lens, cross-task factor disjointness).
%
% v6 changes from v5.5:
%   - NEW CONTRIBUTION C4: Attribution-causal orthogonality + factor manifold
%     enables SOTA DAS. ~88 deg between PCA/Tucker and DAS subspaces.
%     14-40x activation-causal gap. Background/modulation decomposition
%     resolves the gap (Tucker modes necessary but not sufficient).
%   - NEW Section 5.x: Factor overlap across circuit stages (4x4 Jaccard
%     heatmap, within-pathway 0.018-0.037, cross-pathway 0.001-0.004).
%     Per-edge factor fingerprints. Q/K/V factor sharing.
%   - NEW Section 7: Attribution vs Causation Are Orthogonal in Factor Space
%     (threeway subspace comparison, activation-causal gap, background/
%     modulation decomposition, per-counterfactual stability)
%   - NEW Appendix: Per-counterfactual Tucker analysis (V10 experiment A)
%   - NEW Appendix: Distributed logit lens on DAS dimensions
%   - NEW Appendix: Subspace trimming and k-sweep
%   - NEW Appendix: Factor mediation / causal discovery
%   - NEW Appendix: Conceptor steering via factor subspaces
%   - NEW Appendix: Post-hoc sparsity recovery (negative result)
%   - NEW Appendix: Extended null controls
%   - NEW Appendix: Merullo composition comparison
%   - UPDATED abstract to include C4 (orthogonality finding)
%   - UPDATED intro contributions list (4 contributions, not 3)
%   - UPDATED discussion to cover orthogonality interpretation
%   - UPDATED figure/table inventory for new content
%   - DAS/manifold material structured for possible later extraction
%     into standalone paper
%
% Preserved from v5.5:
%   - ALL appendix content, figure specs, table specs, reviewer annotations
%   - ALL technical detail (Tucker modes, EAP derivation, gauge freedom, etc.)
%   - ALL bolded openers, source/data annotations, REVIEWER comments
%   - Model equivalence (Section 3.3)
%   - Family of adapted methods (Section 4) including full DAS section
%   - Full scoring results and ablations
%   - Complete reviewer critique summary
%
% Design principles (from v5.5):
%   - Nanda: 1-3 concrete claims, intro is self-contained, every section must justify its existence
%   - Perplexity: Figure 1 on page 1, every section opens with its conclusion not a question
%   - IOI: contributions list, question-style subsection titles, hypothesis -> experiment -> result
%   - Grokking: "N lines of evidence" roadmap, bolded openers, evidence-type labels
%   - MIB: formal metric definitions with equations, bold finding statements
%   - Advisor: sell readers, proofs in appendix, 8-10 pages for top conference
%
% Narrative arc (v6):
%   (1) We factorize weights as W=SF (architecture, training, equivalence)
%   (2) Substituting W=SF into EAP gives per-factor attribution for free
%   (3) PCA/Tucker on the attribution tensor decomposes circuits into sub-circuits
%   (4) But attribution subspaces are orthogonal to causal subspaces (~88 deg)
%   (5) The factor manifold nonetheless enables SOTA causal variable discovery
%       (91.2% hard-IIA via CPCA-init DAS, gradient attractor)
%   (6) Background/modulation decomposition resolves the gap
%
% Key numbers (v6 additions in CAPS):
%   - CE gap <0.02 nats at 99% sparsity (model equivalence)
%   - Bidirectional DAS IIA >= 0.97 (representational equivalence)
%   - PCA sub-circuits: IOI NMI = 0.600 (p=0.004) vs Wang et al. roles
%   - Tucker = PCA factor directions: cosine 1.000, geodesic ~0
%   - ATTRIBUTION-CAUSAL ANGLE: ~88 DEG ACROSS 5 TASKS
%   - ACTIVATION-CAUSAL GAP: 14-40X (r=0.97-0.99 VS 0.03-0.43)
%   - CPCA-init DAS: 91.2% hard IIA vs 73.5% random (k=1), 100% strict (k=4)
%   - Gradient attractor: rho rises 26% -> 83%
%   - Group lasso DAS: ~81 of 8192 factors sufficient for 98.8% regular IIA
%   - CROSS-PATHWAY FACTOR OVERLAP: JACCARD 0.001-0.004
%   - WITHIN-PATHWAY FACTOR OVERLAP: JACCARD 0.018-0.037
%   - TUCKER MODES NECESSARY (REMOVING 8 DROPS LOGIT DIFF 21%) BUT NOT
%     SUFFICIENT (KEEPING ONLY MODES REVERSES LOGIT DIFF TO -30%)
%   - PER-COUNTERFACTUAL STABILITY: MEAN PRINCIPAL ANGLES 58-78 DEG
%
% Target: TMLR (~25 pages main + unlimited appendix).
% Page budgets in comments. Figures and tables inventoried at bottom.

\documentclass[11pt]{article}
\usepackage[margin=1in]{geometry}
\usepackage{amsmath,amssymb}
\usepackage{graphicx}
\usepackage{booktabs}
\usepackage{textcomp}
\usepackage{hyperref}

% ============================================================
% TITLE OPTIONS (pick one)
% ============================================================
%
% OPTION 1 (CURRENT FAVORITE):
%   "Subspace Circuit Discovery via Sparse Transformer Factorization"
%   - "Subspace Circuit Discovery" = the novel concept (what we find)
%   - "Sparse Transformer Factorization" = the method (how we find it)
%   - "Transformer" not "Weight" because we factorize the whole model
%
% OPTION 2 (safe fallback):
%   "Sparse Transformer Factorization for Sub-Edge Circuit Analysis"
%
% OPTION 3 (memorable):
%   "Opening Up the Edges: Per-Factor Attribution via Sparse Transformer Factorization"

\title{Subspace Circuit Discovery via Sparse Transformer Factorization}
\author{Ivan Vegner, Elliot Tower, N.\ Siddharth, Alex Doumas}

\begin{document}
\maketitle

% ============================================================
% ABSTRACT (~200 words)
% ============================================================

\begin{abstract}

% REVIEWER: "what information flows" is too evocative if you can't name
% the information. Safer: "which factor subspaces carry the signal."
% REVIEWER: "lossless" --- only exact for first-order EAP, approximate
% for EAP-IG. Must be precise here.

[1 sentence: Circuit discovery identifies which edges in a transformer's
computational graph matter for a task, but says nothing about \emph{what}
information flows along each edge --- which subspaces of the residual
stream carry the signal.]

[1 sentence: We decompose each weight matrix as $W = SF$ where $F$ is a
shared factor bank and $S$ is a sparse selector, trained by distillation
to faithfully reproduce a pretrained model (GPT-2 Small, CE gap $<$0.02
nats at 99\% sparsity).]

[1 sentence: This factorization makes existing circuit methods ---
EAP, EAP-IG, EAP-GP, DAS --- automatically produce per-factor
attribution scores: each edge score becomes a sum over per-factor
contributions via a purely algebraic identity, yielding a 3D attribution
tensor of shape (writers $\times$ readers $\times$ factors).]

[1 sentence: PCA on this 3D tensor (flattened to edges $\times$ factors)
automatically decomposes circuits into interpretable sub-circuits --- groups
of edges sharing the same factor subspace --- with auto-$k$ at 90\% variance
requiring no per-task tuning; Tucker decomposition of the same tensor
independently recovers identical factor directions (cosine similarity 1.000,
Grassmannian geodesic distance $\approx$ 0), confirming these sub-circuits
are real structure, not decomposition artifacts.]

[1 sentence: The factor bank also provides a geometric inductive bias
for causal variable localization: constraining PCA to the factor
subspace matches DAS at 98.8\% IIA with zero training, and initializing
DAS from this subspace improves hard-example IIA from 73.5\% to 91.2\%.]

[1 sentence (NEW v6): Analyzing the relationship between attribution
and causal structure reveals a striking finding: the factor subspace
explaining most attribution variance is nearly orthogonal
($\sim$88\textdegree) to the causally-active subspace (DAS), with a
14--40$\times$ gap between activation universality and causal specificity.
Tucker modes capture the modulation component --- necessary for task
performance (removing 8 modes drops logit diff by 21\%) but not
sufficient in isolation (keeping only Tucker modes reverses the logit
diff). The factor manifold acts as a gradient attractor ($\rho$:
26\%$\to$83\%), determining whether DAS training finds the global optimum
--- a launch pad for causal discovery, not a container.]
\end{abstract}


% ============================================================
% 1. INTRODUCTION (~1.5 pages)
% ============================================================

\section{Introduction}

% Open with the problem, not the method (Nanda: "intro should be self-contained")

% REVIEWER: "There is no training-free method" --- NDM (Huang & Hahn 2025)
% does unsupervised subspace decomposition. Must acknowledge and differentiate.
% REVIEWER: "analogous to SAE dictionary" --- reviewer WILL ask how factors
% compare to SAEs. Need to be precise about the differences.

[1 paragraph: Edge-level circuit discovery (EAP, ACDC) tells us \emph{which}
attention heads communicate for a task, but each edge is a single scalar.
We don't know which dimensions of the 768-dimensional residual stream
carry the signal, whether an edge transmits one signal or several
superimposed ones, or whether two edges carry the same information.
Causal variable localization (DAS) finds the relevant subspace but
requires supervised training on counterfactual pairs --- specifying the
causal model in advance. There is no general-purpose method to decompose
an edge into its constituent factor-level contributions.]

% [FIGURE 1 --- NEED to create]
% A single figure that tells the whole story. Three panels:
%
% Panel A: "Standard circuit" --- a small computational graph (4-5 heads)
% with thick/thin edges labeled by EAP scores. Each edge is a single number.
% Caption: "EAP: which edges matter (but what flows along them?)"
%
% Panel B: "Factorized circuit" --- same graph but each edge is now a bundle
% of colored strands (factors). Some strands shared between edges.
% Caption: "Factorized EAP: per-factor attribution for every edge"
%
% Panel C: "Attribution PCA" --- the scores_W matrix (edges x factors) with
% a scree plot inset, then an arrow to a low-rank version. Below: a bar
% showing CMD = 0.023 (PCA k=8) vs 0.052 (full) vs 0.071 (DAS k=4).
% Caption: "SVD denoises the attribution tensor --- fewer dimensions, better circuits"
%
% This must be on page 1.
%
% REVIEWER: Figure 1 Panel C should NOT be PCA-centric if the headline
% is the factorization itself. Consider: Panel C shows the family of
% adapted methods (EAP -> factorized EAP, DAS -> factorized DAS) instead.

\textbf{[FIGURE 1 --- NEED: Three-panel pipeline figure. (A) Standard EAP circuit
with scalar edges. (B) Factorized circuit with per-factor edge bundles.
(C) PCA denoising: scree plot + CMD comparison bar chart.]}

% The approach, compressed

[1 paragraph: We propose \emph{sparse weight factorization}: every weight
matrix $W$ in a pretrained transformer is decomposed as $W = SF$, where
$F$ is a shared overcomplete factor bank and $S$ is a sparse per-projection
selector, trained by distillation on the composite QK and OV circuits.
The model computes through factors --- $Wx = S(Fx)$ --- so factors are
part of the computation, not a post-hoc reconstruction. This structural
property makes existing circuit methods automatically produce per-factor
scores: factorized EAP decomposes each edge score into a sum over
per-factor contributions via a purely algebraic identity. Factorized DAS
identifies which factors span the causal subspace. The contribution is
the factorized model and the adapted methods; PCA sub-circuit
decomposition is the primary tool applied to the factor-level scores,
with Tucker decomposition providing independent validation.]

% The punchline result

% REVIEWER: "The contribution is the factorized model" --- a reviewer
% will say "so you built a model. What did you LEARN from it?"
% The answer: the IOI case study in Section 5.

[1 paragraph: Our main finding is that per-factor attribution reveals
structure invisible to scalar edge scores. PCA on the per-factor
attribution tensor automatically decomposes circuits into interpretable
sub-circuits --- groups of edges sharing the same factor subspace --- and
Tucker decomposition independently recovers identical factor directions
(cosine similarity 1.000, Grassmannian geodesic $\approx$ 0, 8/8
top-edge match on IOI). This convergence of two mathematically different
methods is strong evidence that the sub-circuits are real. PCA at
auto-$k$ (90\% variance) requires no per-task tuning: IOI $k{=}108$,
SVA $k{=}5$, GT $k{=}79$, Gender $k{=}190$, Capital $k{=}29$ ---
adapting to circuit complexity automatically. The right unit of analysis
is the \emph{factor subspace}: individual factor attribution is near
chance (gauge-dependent), but the PCA/Tucker mode spans are
gauge-invariant and causally meaningful. The factor bank also serves as
a geometric inductive bias for causal variable localization: constrained
PCA (restricting to the top 1\% of factors) matches DAS at 98.8\% IIA
with zero training on regular examples, and initializing DAS from this
subspace improves hard-example IIA from 73.5\% to 91.2\% despite
starting 82.5$^\circ$ from the optimum. Factor-constrained subspaces
]

% Contributions (IOI paper style, numbered)

[Contributions:
\begin{enumerate}
\item \textbf{Sparse weight factorization} ($W = SF$): a shared factor
bank + sparse selectors, trained by composite-circuit distillation
(QK, OV), achieving Pareto-optimal sparsity vs.\ fidelity across 7
selector mechanisms and bank sizes from 1024 to 40000. The factorized
model is causally equivalent to the pretrained model (CE gap $<$0.02
nats, DAS IIA $\geq$0.97 bidirectional). (\S\ref{sec:factorization})

\item \textbf{Factorized circuit methods}: algebraically lossless
per-factor decomposition of EAP edge scores, extended to EAP-IG,
EAP-GP, and DAS. Each adapted method returns standard outputs plus
per-factor attribution at zero additional cost.
(\S\ref{sec:factorized_methods}; proof in Appendix~\ref{app:eap_derivation})

\item \textbf{Sub-circuit decomposition via factor subspaces}:
PCA on the per-factor attribution tensor automatically decomposes any
circuit into interpretable sub-circuits, each mediated by a distinct
factor subspace (IOI: early heads $\to$ MLP0, S-inhibition fan-out,
name movers $\to$ logits). Tucker decomposition independently recovers
identical factor directions (cosine 1.000, Grassmannian geodesic
$\approx$ 0), confirming the sub-circuits are real structure, not
decomposition artifacts. The factor \emph{subspace} is gauge-invariant
and causally meaningful; individual factors are not.
(\S\ref{sec:analysis})

\item \textbf{Causal variable localization via the factor manifold}:
Constraining PCA to factor subspace matches DAS on regular examples
(98.8\% IIA, zero training). Manifold-initialized DAS achieves 91.2\%
hard IIA vs 73.5\% vanilla.
% SOURCE: lib/factorized_das/RESULTS_CATALOG.md (IIA numbers)
(\S\ref{sec:factorized_methods})

\item \textbf{Attribution-causal orthogonality and the background/modulation
decomposition} (NEW v6): The factor subspace explaining most attribution
variance is $\sim$88\textdegree\ from the causally-active subspace across
5 tasks, with a 14--40$\times$ activation-causal gap. Tucker C modes
capture the modulation component --- necessary (removing 8 modes drops
logit diff 21\%) but not sufficient (keeping only modes reverses logit
diff to $-$30\%). Different corruptions route through partially distinct
factor subspaces (principal angles 58--78\textdegree). Each circuit
pathway uses nearly disjoint factor subsets (cross-pathway Jaccard
$<$ 0.5\%). (\S\ref{sec:orthogonality})
\end{enumerate}]


% ============================================================
% 2. BACKGROUND (~0.75 pages)
% ============================================================

\section{Background}

% Open with what the reader needs, not a literature survey

% REVIEWER: Must cite Syed et al. 2024, Nanda & Bloom 2022,
% Conmy et al. 2023 (ACDC). Don't miss any.

% --- 2.1 ---
\subsection{Edge Attribution}

[3-4 sentences: Attribution patching approximates activation patching
via first-order Taylor expansion, giving per-module importance in O(1).
EAP (Syed et al.\ 2024) extends this to edges --- the connection from
writer $u$ to reader $v$ through the residual stream.
State the EAP formula:
$\Delta L_{uv} = \sum_n \langle r_u'^{(n)} - r_u^{(n)},\;
\partial L / \partial \text{resid}_v^{(n)} \rangle$.
Each edge gets one scalar: important, but no information about
\emph{what} flows.]

% --- 2.2 ---
\subsection{Causal Abstraction}

[3-4 sentences: Geiger et al.\ (2024) frame interpretability as alignment
between a neural network and a causal model. DAS learns a rotation $R$
that aligns the causal variable with the first $k$ axes, then swaps those
$k$ dimensions between inputs. Interchange Intervention Accuracy (IIA):
fraction where the patched model predicts the source answer. This
identifies \emph{where} information is encoded but requires specifying
the causal model and training $R$.]

% --- 2.3 ---
\subsection{MIB Evaluation Framework}

[3-4 sentences: MIB (Mueller et al.\ 2025) standardizes evaluation.
Track 1 measures circuit faithfulness at multiple sizes via activation
patching; CPR = $\int_0^1 f(\mathcal{C}_k)\,dk$ (area under faithfulness
curve, higher = better) and CMD = $\int_0^1 |1 - f(\mathcal{C}_k)|\,dk$
(distance from perfect reproduction, lower = better).
State the faithfulness formula:
$f(\mathcal{C}, \mathcal{N}; m) = (m(\mathcal{C}) - m(\varnothing)) /
(m(\mathcal{N}) - m(\varnothing))$.]


% ============================================================
% 3. SPARSE WEIGHT FACTORIZATION (~1 page)
% ============================================================

\section{Sparse Weight Factorization}
\label{sec:factorization}

% REVIEWER: "Why not just use SAEs?" This must be addressed HEAD-ON.
% Key differences: (1) factors are in the model's forward pass,
% (2) shared bank across all projections enables cross-layer tracking,
% (3) no reconstruction error --- W = SF exactly (within distillation loss).
% REVIEWER: "You trained a different model and are analyzing THAT, not
% the original. How do you know the factorized model's structure reflects
% GPT-2's structure and not an artifact of the factorization?"
% Answer: model equivalence (Section 3.3) + the factorization is heavily
% constrained (only S and F are free; all other params frozen from GPT-2).

% --- 3.1 Architecture ---
\subsection{Architecture}

% Open with the conclusion (Perplexity guide principle)

[Bolded opener: \textbf{Every weight matrix in the model can be faithfully
decomposed as $W_i = S_i F$, where $F$ is a shared overcomplete factor
bank and $S_i$ is a sparse selector, achieving 99\% sparsity with no
CE degradation.}]

% --- Architecture ---

[1 paragraph: $F \in \mathbb{R}^{N_F \times d_m}$ is a factor bank
whose rows are directions in the residual stream, shared across all
projections ($W_Q, W_K, W_V, W_O, W_{\text{in}}, W_{\text{out}}$).
$S_i \in \mathbb{R}^{d_h \times N_F}$ is a per-projection sparse
selector. The bank is overcomplete ($N_F \geq d_m$; we use $N_F = 8192$
for GPT-2's $d_m = 768$). Each channel in each head activates only a
small subset of factors, so the selector matrix is highly sparse ---
analogous to an SAE's encoder activations, but in \emph{weight space}.
The key difference from SAEs: factors are part of the model's computation
(the model literally computes $W_i x = S_i F x$), not a post-hoc
decomposition of what the model outputs.]

% --- 3.2 Why QK and OV, not individual weights (gauge freedom) ---
\subsection{Distillation on Composite Circuits}

[1 paragraph: We distill on the \emph{composite} circuits
$QK_i = W_Q^i (W_K^i)^\top$ and $OV_i = W_V^i W_O^i$, not on
individual weight matrices. This is necessary because of rotational
gauge freedom: for any orthogonal $R$, replacing $W_Q \to W_Q R$
and $W_K \to W_K R$ leaves the QK circuit unchanged but produces a
completely different individual factorization. The composite circuits
are gauge-invariant observables; matching them is both necessary and
sufficient for faithful reproduction. Formal statement and proof in
Appendix~\ref{app:gauge}.]

% --- Training ---

[2-3 sentences: Core loss: MSE on QK and OV circuits summed over layers.
Sparsity via selector mechanism (DST or JumpReLU with straight-through
estimator). Only $F$ and $S$ optimized; all other parameters copied
from the pretrained model. Full training details in
Appendix~\ref{app:training}.]

% --- 3.3 Model Equivalence --- [NEW IN V4]
\subsection{Model Equivalence}
\label{sec:equivalence}

% REVIEWER: "CE gap is not model equivalence. The model could have
% the same CE but completely different internal representations."
% Answer: we also show CREH L3 (DAS IIA bidirectional) and MAS >= 0.97.
% REVIEWER: "Why always DAS for equivalence? What about other probes?"
% Answer: CE is behavioral. DAS/MAS is representational. EAP recovery
% is structural. Three independent methods at three levels.

[Bolded opener: \textbf{The factorized model is not merely an approximation
of GPT-2 --- it is causally equivalent. CE gap $<$0.02 nats, DAS IIA
$\geq$0.97 bidirectional on IOI, and MAS confirms the factorized
model's representations are interchange-compatible with GPT-2's.}]

[1 paragraph: Three levels of equivalence evidence:
(1) \textbf{Behavioral}: CE gap $<$0.02 nats at 99\% selector sparsity.
    The factorized model produces the same next-token distribution as GPT-2.
(2) \textbf{Representational}: DAS trained on GPT-2 representations achieves
    IIA $\geq$0.97 when applied to the factorized model's representations,
    and vice versa. This bidirectional interchange (MAS) confirms the
    internal representations are structurally aligned, not just
    input-output equivalent.
(3) \textbf{Circuit-level}: EAP-IG on the factorized model recovers the
    same circuit structure as on GPT-2 (same top edges, same head rankings).
    The factorized model eliminates negative name movers --- heads that exist
    in GPT-2 to cancel other heads' errors --- because the factorized
    representation doesn't need that cancellation structure.]

% REVIEWER: "Negative name mover elimination is interesting but you
% buried it in one sentence. This could be a finding."
% Consider: expand in Section 5 or Discussion.

% --- 3.4 Pareto result ---
\subsection{Sparsity--Fidelity Tradeoffs}

% [FIGURE 2: Pareto frontier --- NEED to create or clean up from existing]
% x-axis: average active factors per channel (L0). y-axis: cross-entropy loss.
% GPT-2 baseline as horizontal line. Each point = one checkpoint.
% Color by selector type (DST, JumpReLU, L1). DST dominates.
% Clean, publication-quality version of the W&B scatter.
\textbf{[FIGURE 2 --- NEED: Pareto frontier. CE vs L0/channel by selector
type. DST reaches GPT-2 CE at $\sim$100 active factors/channel (99\% sparse).
Horizontal line at GPT-2 baseline CE.]}

[2 sentences: DST achieves Pareto-optimal tradeoffs: 99\% sparsity
with $<$0.02 CE gap from GPT-2 ($\sim$3.25). The overcomplete factor bank
($N_F = 8192 \gg d_m = 768$) enables extreme sparsity without
information loss.]


% ============================================================
% 4. FACTORIZED CIRCUIT METHODS (~1 page) [NEW SECTION IN V4]
% ============================================================

\section{Factorized Circuit Methods}
\label{sec:factorized_methods}

% REVIEWER: "This is just substituting W=SF into existing formulas.
% Is that really a contribution?"
% Answer: Yes, because (1) the algebraic identity is non-obvious
% (involves LayerNorm Jacobian), (2) it enables a FAMILY of methods,
% not just one, (3) the resulting per-factor scores are a new object
% nobody has had access to before, and (4) the closed-form version
% avoids materializing factor activations entirely.
%
% REVIEWER: "You adapted 4 methods but only show detailed results
% for EAP-IG. What about EAP-GP and DAS results?"
% Must show at least CMD numbers for each adapted method, even briefly.

[Bolded opener: \textbf{The $W = SF$ structure makes four existing
circuit methods --- EAP, EAP-IG, EAP-GP, and DAS --- automatically
produce per-factor attribution scores via algebraic decomposition,
requiring no additional forward passes.}]

% --- 4.1 Factorized EAP ---
\subsection{Factorized EAP and EAP-IG}

[1 paragraph: We first state standard EAP (Eq.~\ref{eqn:eap}): the edge
score $\Delta L_{uv}$ between writer $u$ and reader $v$ is the sum over
tokens of the writer's output difference dotted with the reader's gradient.
Show the formula with shape annotations (as in Ivan's draft \S2.1). Then
substitute $W_u^{o\top} = S_u F$ and $W_v^{in} = F^\top S_v^\top$,
expand, and rearrange to get
$\Delta L_{uv} = \sum_{i,j} \Delta L(u,i,v,j)$ (Eq.~\ref{eqn:feap}).
The per-factor score $\Delta L(u,i,v,j)$ measures how much writing by
module $u$ into factor $i$, as read by module $v$ through factor $j$,
affects the loss. This decomposition is exact --- purely algebraic, no
approximation. We present EAP (not EAP-IG) because the linear
approximation makes the algebra exact; EAP-IG uses integrated gradients
and the decomposition becomes approximate.
Full derivation with LayerNorm Jacobian expansion in
Appendix~\ref{app:eap_derivation}.]

% --- Losslessness ---

[1 paragraph: State the losslessness formally as a proposition.
\textbf{Proposition 1} (Factorized EAP aggregates to EAP):
$\Delta L_{uv} = \sum_{i=0}^{N_F} \sum_{j=0}^{N_F} \Delta L(u,i,v,j)$.
This is immediate from the derivation --- factorized EAP is strictly
more informative than standard EAP at zero additional cost (one forward
+ one backward pass, same as standard EAP). The closed-form score
$\Delta L(u,i,v,j) = M_{ij} T_1^{(uv)} - \frac{1}{d_m} T_2^{(uv)}$
(where $M$ is the factor Gram matrix and $T_1, T_2$ are sum-of-outer-products)
enables efficient computation via einsums.
Closed-form derivation in Appendix~\ref{app:closed_form}.]

% --- 4.2 Factorized EAP-GP --- [NEW IN V4]
\subsection{Factorized EAP-GP (GradPath)}

[1 paragraph: EAP-GP (gradient path attribution) tracks gradients
through the full computational graph. The $W = SF$ substitution
applies identically: each GradPath edge score decomposes into
per-factor contributions. Brief comparison of EAP-GP CMD vs
EAP-IG CMD on IOI.]

% --- 4.3 Factorized DAS --- [NEW IN V4, EXPANDED IN V5, REWRITTEN V5.5]
\subsection{Factorized DAS and Causal Subspace Localization}
\label{sec:factorized-das}

% REVIEWER: "DAS on the raw model also finds a subspace. What's the
% advantage of constraining to factor span?"
% Answer: (1) Training-free constrained PCA matches DAS on regular
% examples (98.8% IIA). (2) Factor-manifold initialization dramatically
% improves DAS on hard examples (91.2% vs 73.5%). (3) Extensive ablation
% shows that training is required for hard examples --- the manifold is
% an initialization aid, not a substitute.

% SOURCE: writeups/SUB_EDGE_RESULTS.md lines 73-82 (IIA 0.988 at k=2..16)
% SOURCE: gradient_and_geodesic_preview_v2.tex (CPCA-init 91.2%, random-init 73.5%, strict IIA 100% at k=4)
% SOURCE: artifacts/geodesic_line_search_5bases/ (5-bases geodesic results)
% SOURCE: Modal hard-init-methods job (8 untrained methods results)
% DATA: factorized_das results .pt files (per_proj_sel format)
% DATA: artifacts/geodesic_line_search_5bases/geodesic_line_search_5bases_k1.json
% DATA: artifacts/geodesic_line_search_5bases/geodesic_line_search_5bases_k4.json
% DATA: Modal fc-results volume /hard_init_methods_k1/, /hard_init_methods_k4/

\textbf{Standard IIA on regular examples is a deceptively easy metric:
zero-training constrained PCA achieves 98.8\%, trained DAS achieves
94--97\%, and even random baselines exceed 90\% because fewer than 2\%
of examples are flippable.}
Hard examples (logit margin $< 1.0$) are the only regime that
differentiates methods.
On hard examples, CPCA-initialized DAS achieves the best results of
any initialization strategy (91.2\% at $k=1$, 97.1\% at $k=4$, and
100\% strict IIA at $k=4$), while random-init DAS reaches only 73.5\%.
Thirteen untrained methods and five structured-basis geodesic searches
all fail to exceed the majority-class baseline (70.6\%) on hard examples,
confirming that training is required --- but the factor manifold
determines \emph{where} on the Grassmannian that training succeeds.

\subsubsection{Regular Examples: CPCA Matches Trained DAS}

% SOURCE: writeups/SUB_EDGE_RESULTS.md lines 73-82 (CPCA 98.8% IIA)
% SOURCE: lib/factorized_das/RESULTS_CATALOG.md (DAS IIA numbers, layer choices)
Standard DAS learns a rotation $R$ in $\mathbb{R}^{d_m}$.
Factorized DAS constrains $R$ to the span of the factor bank via
$U = F^\top A$ where $A \in \mathbb{R}^{N_F \times k}$ is sparse
(group lasso penalty). But a simpler approach works for regular examples:
\emph{constrained PCA} (CPCA) selects the top $\sim$1\% of factors by selector
magnitude ($\sim$81 of 8192), projects activation deltas into this
81-dimensional subspace, and takes the top-$k$ principal components.
This zero-training method achieves 98.8\% IIA on IOI at Layer~9 ---
compared to 21.3\% for unconstrained PCA at Layer~8. The factor bank
collapses the search space from $\mathrm{Gr}(k, 768)$ to
$\mathrm{Gr}(k, 81)$, making PCA sufficient where it normally fails.

This near-perfect regular IIA is \emph{not} evidence that the causal
variable is easy to find --- it reflects the metric's insensitivity.
On the regular evaluation set, only 2/370 IOI and 5/380 SVA examples
are flippable, so any reasonable direction achieves $\geq$99\% IIA.
Hard examples are where methods actually separate.

\subsubsection{Hard Examples: CPCA-Init DAS is the Best Strategy}

% SOURCE: gradient_and_geodesic_preview_v2.tex Section 1.2 (CPCA-init 91.2%, random 73.5%)
% SOURCE: gradient_and_geodesic_preview_v2.tex line 56 (strict IIA 100%/100% at k=4)
% SOURCE: lib/factorized_das/modal_riemannian_das_experiment.py (Riemannian DAS 64%)
% DATA: Modal fc-results volume /riemannian_das/ (2x2 design results)
On hard examples (logit margin $< 1.0$, $\sim$8\% of data),
CPCA alone achieves only 47--53\% IIA. But
\emph{initializing} vanilla DAS from the CPCA direction ---
then training in unconstrained $\mathbb{R}^{768}$ --- yields the best
performance of any initialization strategy:
91.2\% hard IIA at $k=1$ (100 steps), versus 73.5\% for random-init DAS.
At $k=4$, CPCA-init DAS achieves 97.1\% IIA and \textbf{100\% strict IIA}
on both IOI and SVA, while unconstrained PCA init achieves only
66.7\%/90\% strict IIA.
The CPCA direction is 82.5$^\circ$ from the DAS
optimal (nearly orthogonal on $\mathrm{Gr}(1, 768)$) --- the factor
manifold provides a better \emph{optimization basin}, not a better
initial direction.
More constraint is worse: Riemannian DAS (Stiefel manifold constraint)
drops to 64\% on regular IOI. The factor manifold is a launch pad, not
a prison.

\begin{table}[h!]
\centering
\small
\caption{DAS IIA on hard IOI examples (34 eval, Layer~8). CPCA-init unconstrained achieves the best result across all four conditions.}
\label{tab:das-2x2}
\begin{tabular}{lcc}
\hline
 & \textbf{Unconstrained} & \textbf{On-manifold (Riemannian)} \\
\hline
\textbf{Random init}   & 73.5\% & 88.2\% \\
\textbf{CPCA init}     & \textbf{91.2\%} & 64.0\% \\
\hline
\end{tabular}
\end{table}

\subsubsection{Riemannian Gradient Decomposition}

% SOURCE: gradient_and_geodesic_preview_v2.tex Section 1.1 (rho rises 26%->82.5%)
% SOURCE: gradient_and_geodesic_preview_v2.tex Section 1.2 (geodesic interpolation, angular distances)
% DATA: Modal fc-results volume /gradient_decomposition/ (per-step rho, angular trajectories)
Why does CPCA initialization outperform all alternatives?
We decompose each DAS gradient step into components tangent and normal
to the factor manifold $\mathrm{Gr}(1, 81) \subset \mathrm{Gr}(1, 768)$.
The manifold gradient fraction $\rho = \|\nabla_{\text{on}}\| / \|\nabla_{\text{tang}}\|$
measures how much useful gradient signal lies along the factor manifold.

Two findings explain CPCA-init's superiority:
(1)~$\rho$ \emph{rises} from 26\% to 82.5\% for CPCA init, while
staying flat at $\sim$50\% for unconstrained PCA init ---
the factor manifold is a \emph{gradient attractor} (Figure~\ref{fig:grassmannian-landscape}).
(2)~CPCA init converges to the global optimum (0.1$^\circ$ away);
unconstrained PCA init converges to a local optimum 23.8$^\circ$ away,
despite starting 17$^\circ$ closer.
On $\mathrm{Gr}(4, 768)$, the effect is stronger: CPCA init converges to
$0.59^\circ$ from optimal while unconstrained PCA gets stuck at $86.5^\circ$.

\begin{figure}[h!]
\centering
\includegraphics[width=0.48\textwidth]{figures/grassmannian_landscape/grassmannian_landscape_v11_dark_optimal.png}
\hfill
\includegraphics[width=0.48\textwidth]{figures/grassmannian_landscape/grassmannian_landscape_v12_rho.png}
\caption{\textbf{Grassmannian landscape for DAS training on hard IOI, $\mathrm{Gr}(1, 768)$.}
\emph{Left:} IIA as a function of angle to the DAS optimal direction ($x$-axis) and
factor manifold proximity ($y$-axis). The on-manifold init (blue, starting at 100\%
proximity) follows a high-IIA ridge to the global optimum (star); the off-manifold
init (orange, starting at $\sim$12\% proximity) gets trapped in a local optimum.
\emph{Right:} Same landscape, $x$-axis replaced by manifold gradient fraction $\rho$.
The on-manifold trajectory's $\rho$ rises from 0.26 to 0.83 (gradient attractor);
the off-manifold trajectory stays flat at $\sim$0.35.}
\label{fig:grassmannian-landscape}
\end{figure}

\subsubsection{Geodesic Line Search: Training-Free Methods Fail on Hard Examples}

% SOURCE: lib/factorized_das/geodesic_line_search_5bases.py (experiment code)
% SOURCE: lib/factorized_das/modal_geodesic_line_search_5bases.py (Modal wrapper)
% DATA: artifacts/geodesic_line_search_5bases/geodesic_line_search_5bases_k1.json
% DATA: artifacts/geodesic_line_search_5bases/geodesic_line_search_5bases_k4.json
Can we find high-IIA directions without gradient descent?
We implement a geodesic line search: Riemannian gradient at $Q_0$,
30-point coarse sweep along the geodesic, 10-step ternary refinement
($\sim$40 evaluations total, vs.\ 200 gradient steps for DAS).
We compare five initialization bases --- factor bank (CPCA), SAE decoders,
unconstrained PCA, probing directions, and random --- each with 10--11 starts.

\begin{table}[h!]
\centering
\small
\caption{Geodesic line search IIA on hard IOI (34 eval, Layer~8). Mean $\pm$ std over 10--11 runs. The 70.6\% floor ($= 24/34$) is the majority-class baseline.}
\label{tab:geodesic-5bases}
\begin{tabular}{lcccc}
\hline
\textbf{Basis} & \textbf{Mean ($k$=1)} & \textbf{Max} & \textbf{Mean ($k$=4)} & \textbf{Max} \\
\hline
Probing             & 81.3 $\pm$ 2.0 & 82.4 & 80.5 $\pm$ 3.0 & 82.4 \\
Unconstrained PCA   & 72.7 $\pm$ 7.1 & 94.1 & 73.3 $\pm$ 7.9 & 97.1 \\
Factor bank (CPCA)  & 70.6 $\pm$ 0.0 & 70.6 & 70.9 $\pm$ 0.9 & 73.5 \\
SAE decoders        & 70.6 $\pm$ 0.0 & 70.6 & 71.1 $\pm$ 1.2 & 73.5 \\
Random              & 70.9 $\pm$ 0.9 & 73.5 & 70.6 $\pm$ 0.0 & 70.6 \\
\hline
\textit{Trained DAS (CPCA init)}    & \multicolumn{2}{c}{\textit{94.1}} & \multicolumn{2}{c}{\textit{97.1}} \\
\hline
\end{tabular}
\end{table}

A single geodesic step is insufficient: nearly all bases converge to
the majority-class baseline (70.6\% $= 24/34$). Probing, which is
supervised, reaches 81\% but cannot flip the hard examples.
The factor bank (CPCA) is trapped at the baseline from every start ---
the geodesic cannot escape the local landscape.
This is why CPCA \emph{initialization} for trained DAS works while
CPCA alone does not: the factor manifold identifies the right basin,
but reaching the optimum within that basin requires gradient descent
through the IIA objective.

\subsubsection{Comprehensive Untrained-Method Comparison}

% SOURCE: lib/factorized_das/hard_init_methods.py (experiment code: 8 methods)
% SOURCE: lib/factorized_das/modal_hard_init_methods.py (Modal wrapper)
% DATA: Modal fc-results volume /hard_init_methods_k1/, /hard_init_methods_k4/
We further test 8 zero-training or low-training methods:
CPCA baseline, difference-in-means, top EAP factors (gradient-weighted
attribution), logit gradient (1 backward pass), constrained probing
(200 steps in factor subspace), LDA in factor subspace, hard-example
CPCA (PCA on hard deltas only), margin-weighted CPCA, and iterative
geodesic descent (5 Riemannian steps from CPCA).

\begin{table}[h!]
\centering
\small
\caption{Hard-example IIA for untrained methods vs.\ trained DAS.
All factor-constrained methods have manifold fraction $\approx 1.0$.
No untrained method exceeds the majority-class baseline (70.6\%).}
\label{tab:untrained-methods}
\begin{tabular}{lcc}
\hline
\textbf{Method} & \textbf{Hard IIA ($k$=1)} & \textbf{Hard IIA ($k$=4)} \\
\hline
Iterative geodesic      & 70.6 & 70.6 \\
LDA (factor subspace)   & 70.6 & 67.6 \\
Constrained probing     & 67.6 & 64.7 \\
Top EAP factors         & 64.7 & 47.1 \\
Logit gradient          & 50.0 & 52.9 \\
Margin-weighted CPCA    & 52.9 & 50.0 \\
CPCA baseline           & 52.9 & 47.1 \\
Hard-example CPCA       & 52.9 & 47.1 \\
Difference-in-means     & 52.9 & 47.1 \\
\hline
Trained DAS (CPCA init) & \textbf{91.2} & \textbf{97.1} \\
Trained DAS (random)    & 73.5 & --- \\
\hline
\end{tabular}
\end{table}

No untrained method exceeds the majority-class baseline (70.6\%).
The 70.6\% barrier corresponds to 24/34 examples that are already
correctly predicted; the remaining 10 require the intervention to
\emph{flip} the model's prediction. At 70.6\%, none of these 10 are
flipped. Only trained DAS breaks through, and CPCA-init DAS breaks
through the farthest.

Hard-example CPCA (training PCA only on hard deltas) provides no
benefit over standard CPCA --- the hard examples do not occupy a
spectrally distinct subspace. They differ in \emph{degree} (smaller
logit margin), not \emph{kind} (different direction).

\paragraph{Takeaway.}
Regular IIA is not the right metric for evaluating causal variable
discovery --- it is too easy, saturating at $\geq$94\% for any
reasonable method. Hard examples are where methods differentiate, and
on hard examples CPCA-initialized DAS is the best strategy we tested:
91.2\% at $k=1$, 97.1\% at $k=4$, and 100\% strict IIA at $k=4$.
The factor manifold provides a $20\times$ search space reduction
($\mathrm{Gr}(k, 81)$ vs.\ $\mathrm{Gr}(k, 768)$), places
initialization in the global optimum's basin of attraction
(Figure~\ref{fig:grassmannian-landscape}), and acts as a gradient
attractor during training ($\rho$ rises from 26\% to 83\%).
It cannot replace training --- but it determines whether training
finds the right answer.

% --- 4.4 The attribution tensor ---
\subsection{The Per-Factor Attribution Tensor}

[1 paragraph: For circuit ranking, we marginalize over the reader
factor to get a matrix $\text{scores}_W \in \mathbb{R}^{n_{\text{edges}}
\times N_F}$ (69,865 $\times$ 8,192 for GPT-2 IOI). Each row is an
edge's factor-attribution profile. Standard EAP is the row-sum
(collapsing factor dimension). Factorized EAP preserves the full
factor profile per edge. Or keep the 3D tensor
$\mathbb{R}^{n_w \times n_r \times N_F}$ for Tucker decomposition.
The question is: what is the intrinsic dimensionality of this
factor-attribution space?]


% ============================================================
% 5. FACTOR-LEVEL CIRCUIT ANALYSIS: IOI CASE STUDY (~2 pages)
% ============================================================

\section{Factor-Level Circuit Analysis: IOI Case Study}
\label{sec:analysis}

% REVIEWER: "This is all on one task (IOI). How do you know it generalizes?"
% Answer: concentration analysis is on 3 tasks. Cross-architecture is
% in Section 6. But IOI is the worked example.
% REVIEWER: "You're analyzing the factorized model, not GPT-2.
% Everything you find could be an artifact of the factorization."
% Answer: Section 3.3 establishes causal equivalence. The factorized
% model preserves GPT-2's circuits (same edges, same rankings).
% The factor structure is HOW the preserved circuits are implemented
% in the factorized representation.

[1 paragraph: We demonstrate what per-factor attribution reveals using
IOI as the primary worked example, with PCA decomposition across 5 tasks
(IOI, SVA, Greater-Than, capital-country, gender bias). PCA on the
flattened attribution tensor automatically decomposes circuits into
sub-circuits: groups of edges sharing the same factor subspace. Each
principal component defines a ``factor direction'' --- edges loading on
the same PC use the same factor subspace to communicate, forming a
coherent sub-circuit. Tucker decomposition of the full 3D tensor
independently recovers identical factor directions (cosine similarity
1.000, 8/8 top-edge match), confirming this structure is real. All use
checkpoint atomic-sweep-40 (8192 factors, DST selector).]

% --- 5.1 PCA Sub-Circuit Decomposition ---
\subsection{PCA Sub-Circuit Decomposition}

[Bolded opener: \textbf{PCA on the flattened attribution tensor
$\text{scores}_W \in \mathbb{R}^{(n_w \cdot n_r) \times N_F}$
automatically decomposes any circuit into interpretable sub-circuits,
each mediated by a distinct factor subspace, with auto-$k$ at 90\%
explained variance requiring no per-task tuning.}]

% SOURCE: Auto-k values computed in pca_multi_task_subcircuits.py (90% cumulative variance threshold)
% DATA: /results/das_eap/atomic-sweep-40/multi_task_pca/{task}/scores_WR_ig10.pt
% --- The tensor and why PCA works ---

[1 paragraph: The factorized attribution tensor
$\text{scores}_W \in \mathbb{R}^{n_w \times n_r \times N_F}$ assigns
each edge a vector in $\mathbb{R}^{N_F}$ rather than a scalar. Flattening
to $\mathbb{R}^{(n_w \cdot n_r) \times N_F}$ and applying truncated SVD
yields principal components --- orthogonal directions in factor space ---
and per-edge loadings on each component. Edges loading on the same PC
share a factor subspace: they communicate through the same subset of
factors and form a coherent sub-circuit. The number of components $k$ is
set automatically at the 90\% variance threshold: IOI $k{=}108$,
SVA $k{=}5$, GT $k{=}79$, Gender $k{=}190$, Capital $k{=}29$. This
requires no per-task tuning and adapts to circuit complexity ---
structurally simple circuits (SVA) concentrate variance in few
components while complex circuits (Gender, IOI) spread it across many.]

% --- Tucker equivalence ---
\subsection{Tucker Decomposition Recovers Identical Factor Directions}
\label{sec:tucker_equivalence}

[Bolded opener: \textbf{Tucker decomposition of the full 3D tensor
$\text{scores}_W \approx \mathcal{G} \times_1 U_w \times_2 U_r \times_3 U_f$
recovers factor modes $U_f$ that are identical to the PCA principal
components (cosine similarity 1.000, Grassmannian geodesic distance
$\approx$ 0, 8/8 top-edge match), providing independent validation that
the sub-circuit structure is real.}]

[1 paragraph: Tucker decomposition preserves the 3D structure of the
attribution tensor, yielding writer modes $U_w$, reader modes $U_r$,
factor modes $U_f$, and a core tensor $\mathcal{G}$. The factor modes
$U_f \in \mathbb{R}^{N_F \times r_3}$ identify directions in factor
space --- the same object PCA finds from the flattened matrix. On IOI
with Tucker ranks $(10,10,8)$: each of the 8 Tucker factor modes aligns
with a PCA PC at $|\cos\theta| \geq 0.998$; the cross-method cosine
similarity matrix is a near-perfect identity matrix (all off-diagonal
entries $< 0.04$); and every matched pair identifies the same dominant
edge. This convergence of two mathematically different methods ---
PCA maximizes variance of the flattened matrix, Tucker minimizes
reconstruction error of the 3D tensor --- is strong evidence that the
factor sub-circuits reflect genuine structure in the attribution tensor,
not artifacts of the decomposition. Tucker's additional writer and reader
modes are fully determined by the (shared) factor modes: each factor
direction is associated with a single dominant writer-reader pattern,
so the extra dimensions add no information beyond what the PCA edge
loadings already contain.]

% SOURCE: pca_tucker_agreement computed in modal_all_figures.py
% DATA: pca_tucker_agreement.json (cosine_matrix, geodesic_matrix)
% VERIFIED: Previous session output showed cos=1.000 on all 8 matched pairs
% [FIGURE --- Agreement matrix]
\textbf{[FIGURE --- PCA vs Tucker factor direction agreement.
8$\times$8 $|\cos|$ matrix = perfect diagonal. Geodesic distance
matrix = all zeros on diagonal, $\pi/2$ off-diagonal. One compact
panel suffices; no side-by-side comparison needed since the methods
are identical. See pca\_tucker\_geodesic\_agreement.png.]}

% --- Method and rank selection ---

\subsubsection{Method and Rank Selection}

[1 paragraph: We compute Tucker decomposition via Higher-Order SVD (HOSVD)
initialization followed by Alternating Least Squares (ALS) refinement,
using tensorly. Ranks $(r_1, r_2, r_3)$ are selected by two criteria:
(1) \textbf{variance elbow}: explained variance saturates at $(10, 10, 8)$
--- increasing to $(15, 15, 8)$ adds $<$2\% variance, while $(5, 5, 4)$
loses 20\% (Figure~\ref{fig:rank_sweep}); (2) \textbf{core tensor
concentration}: at $(10, 10, 8)$, only 5 of 800 core entries account for
50\% of total energy, and 25 entries for 90\%. Higher ranks add near-zero
entries without improving the decomposition. The core tensor's extreme
sparsity is the rank justification: the circuit is dominated by a handful
of writer-reader-factor interactions, and the decomposition has captured
all of them.]

% [FIGURE 3: Rank selection]
\textbf{[FIGURE 3 --- NEED: Two-panel rank justification.
Left: Variance explained vs Tucker rank config ---
(5,5,4), (8,8,6), (10,10,8), (12,12,8), (15,15,8), (20,20,16).
Elbow at (10,10,8).
Right: Cumulative core tensor energy vs number of entries (sorted by
magnitude). At (10,10,8): 5 entries = 50\%, 25 = 90\%, 800 total.
Shows extreme sparsity.]}

\subsubsection{A Priori Rank Selection via Effective Rank}

[Bolded opener: \textbf{The effective rank per mode --- computed by SVD
on each unfolding of the attribution tensor before Tucker --- predicts
both the optimal Tucker configuration and whether Tucker will improve
circuit scoring.}]

[1 paragraph: Before running Tucker, we compute the effective rank of
each mode by SVD on the mode-$n$ unfolding, measuring the number of
singular values needed for 90\% of the energy. This gives an a priori
diagnostic:]

\textbf{[TABLE 1a --- Effective rank per mode across tasks.]}
% Mode             | IOI | SVA | GT  | Hypernymy | Cap-Country | Gender
% Writers  (90%)   | 13  | 5   | 10  | TBD       | TBD         | TBD
% Readers  (90%)   | 17  | 1   | 10  | TBD       | TBD         | TBD
% Factors  (90%)   | 51  | 5   | 51  | TBD       | TBD         | TBD
% s1/s2 ratio (rd) | 2.5 | 8.0 | 5.8 | TBD       | TBD         | TBD

[1 paragraph: SVA's reader effective rank is 1 --- a single reader mode
(m0, MLP layer 0) captures 90\% of the reader-mode energy, with an
$\sigma_1/\sigma_2$ ratio of 8.0. This means the reader dimension is
trivial: all circuit edges route through a single downstream target,
and Tucker's reader decomposition reduces to a scalar multiplier that
cannot change edge rankings. IOI (reader rank 17) and GT (reader rank 10)
have multi-dimensional reader structure that Tucker can exploit. The
effective rank table is both a rank-selection guide (set $r_n$ at or
above the 90\% effective rank) and a diagnostic of circuit complexity:
rank-1 modes indicate structurally simple circuits where the corresponding
dimension offers no denoising opportunity.]

\subsubsection{Scoring Results Across Tasks}

[Bolded opener: \textbf{PCA at auto-$k$ (90\% variance) improves or
matches EAP-IG CMD on 4/5 tasks with zero tuning. Tucker matches PCA's
factor directions exactly but requires per-task rank selection to achieve
its best CMD, and adaptive rank selection (variance-threshold) fails
catastrophically.}]

% SOURCE: tucker_deep_dive/TUCKER_ANALYSIS_COMPLETE.md Section 2 "Results"
% DATA: das_eap/atomic-sweep-40/tucker_scoring/{task}/tucker_scoring.json
% PCA CMD numbers from modal_pca_scoring_variants.py runs
% [TABLE 1b: Scoring results]
\textbf{[TABLE 1b --- PCA and Tucker scoring vs standard EAP-IG.]}
% Task    | EAP-IG CMD | PCA CMD (auto-k@90%) | Best Tucker CMD | Tucker config
% IOI     | 0.032      | 0.021 (k=128)  +34%  | 0.013 (10,10,8)   +59%
% SVA     | 0.015      | 0.014 (k=512)  +5%   | 0.022 (worse)
% GT      | 0.030      | 0.018 (k=256)  +40%  | 0.018 (20,20,16)  +39%
% Gender  | 0.017      | 0.013 (k=512)  +27%  | 0.017 (20,20,16)  +2%
% Capital | 0.021      | 0.025 (k=64)   worse  | 0.021 (no improvement)

[1 paragraph: PCA with auto-$k$ at 90\% variance is the universal
scoring method: it improves CMD on IOI (+34\%), GT (+40\%), Gender
(+27\%), and SVA (+5\%) with no per-task tuning. Tucker at its
best fixed ranks matches or exceeds PCA on IOI (+59\%) and GT (+39\%),
but requires per-task rank selection --- there is no universal Tucker
configuration. Critically, adaptive Tucker (setting ranks to mode
effective ranks at 50--90\% variance) is catastrophically worse than
EAP-IG (CMD 10--27$\times$ higher), because Tucker's scoring value
comes from aggressive compression/denoising, not faithful reconstruction.
PCA operates at 90\% variance (gentle denoising); Tucker must operate
at $\sim$50--60\% variance (aggressive denoising) to improve scoring,
but the optimal aggressiveness varies by task. PCA's simplicity and
universality make it the recommended default; Tucker's value is as an
interpretability tool (\S\ref{sec:tucker_equivalence}), not a scoring
method.]

% [TABLE 1c: Combined Tucker results across tasks]
\textbf{[TABLE 1c --- NEED: Combined table with effective rank, optimal
Tucker config, CMD, core tensor sparsity across all 5 tasks.]}
% This is the main results table for the paper.

% --- Interpretable modes ---

\subsubsection{Interpretable Sub-Circuit Modes}

% SOURCE: tucker_deep_dive/FINDINGS.md (IOI sub-circuit table, lines 37-52)
% DATA: tucker_deep_dive/sub_circuits_ioi.png (top 10 sub-circuits by core energy)
% CROSS-REF: Tucker factor modes match PCA PCs at cos=1.000 (pca_tucker_agreement.json)
[Bolded opener: \textbf{PCA automatically decomposes IOI into three
functional groups --- corresponding to known sub-circuit components ---
without any supervision, circuit labels, or causal model specification.}]

[1 paragraph: The 8 principal components of the IOI attribution tensor
separate into three clean functional groups:
\begin{itemize}
\item \textbf{PC0--PC3 (47\% variance):} Four L0 attention heads
(a0.h1, a0.h4, a0.h5, a0.h3) writing to MLP layer 0 (m0). Each head
occupies its own PC --- they use \emph{different factor subspaces} to
write to the same reader, which is why they separate into 4 components.
These are the duplicate-token and early-processing heads from Wang et al.\ (2023).
\item \textbf{PC4--PC6 (11\% variance):} a5.h5 (the S-inhibition head)
fanning out to mid-layer readers: m5, a6.h3.Q, m6, a7.h0.K, a7.h11.Q,
a8.h10.Q, a8.h11.K. Different PCs capture different reader targets of
the same writer, with different factor subspaces mediating each pathway.
\item \textbf{PC7 (2\% variance):} Name mover heads (a9.h9, a9.h6,
a10.h7, a11.h10) writing to logits. The final output circuit.
\end{itemize}
This is the core finding: \textbf{IOI is not one circuit --- it is a
composition of several sub-circuits, each mediated by a distinct factor
subspace.} The factor bank provides the decomposition; PCA merely
reveals it.

% NOTE: The clustering is fully unsupervised (PCA on the attribution
% tensor, no labels). The role NAMES ("duplicate-token", "S-inhibition",
% "name mover") come from matching discovered head groups to Wang et al.'s
% prior annotations. Frame carefully: "PCA discovers head groupings that
% correspond to known circuit roles" --- not "PCA discovers the roles."
% The discovery is the grouping + factor subspace separation; the naming
% is post-hoc identification with existing labels.]

% [FIGURE 4: Layer-aggregated sub-circuit heatmap]
\textbf{[FIGURE 4 --- HAVE: Layer-aggregated PCA sub-circuit heatmap.
8 panels, one per PC. Each panel is a 13$\times$13 (writer layer $\times$
reader layer) heatmap showing where that PC's energy is concentrated.
PC0--3 light up embed/L0 $\to$ L0. PC4--6 light up L5 $\to$ L5/L6/L7.
PC7 lights up L9/L10/L11 $\to$ logits. Clean separation of the three
functional groups. See pca\_layer\_heatmap.png.

Also: zoomed heatmap showing top 15 edges per PC with head names and
numeric loadings. See pca\_subcircuit\_zoomed.png.

Also: per-task versions with 8 PCs each, showing that different tasks
have different layer structure. See pca\_layer\_heatmap\_\{task\}.png.]}

[1 paragraph: The layer-aggregated view (Figure~4) collapses the
157$\times$445 edge tensor to 13$\times$13 (writer layer $\times$
reader layer) per PC, making the sub-circuit structure immediately
visible. Each PC lights up a different layer pair: PC0--3 concentrate
in embed/L0 $\to$ L0 (early processing), PC4--6 in L5 $\to$ L5--L8
(S-inhibition fan-out), PC7 in L9--L11 $\to$ logits (output). This
decomposition requires no supervised labels, no causal model
specification, and no additional forward passes --- only SVD on the
attribution matrix.]

% --- NEW v6: Factor overlap across circuit stages ---

\subsubsection{Each Edge Uses a Distinct Factor Subset}
\label{sec:factor_overlap}

% SOURCE: today's session (June 19 2026); TUCKER_CIRCUIT_SECTION_FOR_PAPER.md
% DATA: fig3_factor_overlap_stages_v2.png, fig1_per_edge_factor_fingerprints.png,
%       fig2_qkv_factor_sharing.png (all in BATCH2/)

[Bolded opener: \textbf{Circuit stages use nearly disjoint factor subsets.
Even edges within the same pathway share only 1.8--3.7\% of their top-20
factors, and cross-pathway sharing drops to $<$0.5\%.}]

[1 paragraph: For each of the four canonical IOI pathways
(Dup$\to$Ind, Prev$\to$Ind, Ind$\to$SInh, SInh$\to$NM), we pool the
top-20 factors per edge (by attribution magnitude) and compute pairwise
Jaccard similarity. Within-pathway Jaccard ranges from 0.018
(Dup$\to$Ind) to 0.037 (Ind$\to$SInh). Cross-pathway Jaccard is
0.001--0.004. Each edge operates through its own $\sim$20 factors out
of the 8192-dimensional bank. This is not cherry-picking: the pattern
holds across all 6 pathway pairs. Even edges doing the \emph{same job}
(e.g., two S-inhibition$\to$name-mover connections) use almost entirely
different factors --- each edge has a unique factor fingerprint.]

% [FIGURE: Factor overlap heatmap]
\textbf{[FIGURE --- HAVE: 4$\times$4 factor overlap heatmap across IOI
circuit stages. Diagonal (within-pathway) = 0.018--0.037. Off-diagonal
(cross-pathway) = 0.001--0.004. YlOrRd colormap, values 0--0.04.
See fig3\_factor\_overlap\_stages\_v2.png in BATCH2/.
Caption should explain: diagonal values are NOT self-similarity (which
would be 1.0) but mean pairwise Jaccard between different edges within
the same stage.]}

\subsubsection{Per-Edge Factor Fingerprints and Q/K/V Factor Sharing}
\label{sec:factor_fingerprints}

% SOURCE: fig1_per_edge_factor_fingerprints.png, fig2_qkv_factor_sharing.png

[1 paragraph: Per-edge factor fingerprints confirm the disjointness
finding at individual edge resolution: the top-20 factors for each
Ind$\to$SInh edge are largely non-overlapping, with different edges
using different factor channels to carry information between the same
head pairs. Q/K/V projections within a single head use partially
overlapping but not identical factor sets --- Q and K share more
factors with each other than with V, consistent with the $QK^\top$
attention-score duality observed in the 4D Tucker decomposition
(\S\ref{app:tucker_4d}).]

% [FIGURE: Per-edge fingerprints + Q/K/V sharing]
\textbf{[FIGURE --- HAVE: Left: Per-edge factor fingerprints for the 8
Ind$\to$SInh edges (top-20 factors each, binary heatmap showing which
factors each edge uses). Right: Q/K/V factor sharing for 4 IOI heads
with overlap stats.
See fig1\_per\_edge\_factor\_fingerprints.png and
fig2\_qkv\_factor\_sharing.png in BATCH2/.]}

\subsubsection{Cross-Task Factor Disjointness}
\label{sec:cross_task_factors}

% SOURCE: SUB_EDGE_DECOMPOSITION_V4.md

[1 paragraph: At k=10, 76\% of the most-attributed factors are
task-disjoint (38 of 50 unique across 5 tasks). Per-edge Jaccard is
0.007 at k=10, with 89\% of edges having zero factor overlap between
task pairs. Tasks use the same factor bank but activate different
subsets. This is consistent with the cross-pathway disjointness within
IOI: the factor bank provides a large shared vocabulary, but each
circuit edge --- within or across tasks --- selects its own sparse
subset.]

% --- Tucker as interpretability tool on pre-identified circuits ---

\subsubsection{Tucker on Pre-Identified Circuits}

[1 paragraph: Tucker can also be applied to the \emph{circuit sub-tensor}
--- the top-$k$\% edges pre-identified by standard EAP-IG. On the top-1\%
sub-tensor (698 edges), Tucker (10,10,8) explains 93\% variance on IOI,
95\% on SVA, 79\% on GT. The core tensor is sparse on all three tasks:
3--5 entries for 50\% energy. Writer and reader modes are near-axis-aligned
(loading $>$0.99 on individual heads) across all tasks. \emph{This confirms
that Tucker's modular decomposition is a general property of circuit
sub-tensors, not an IOI artifact.} However, the interpretable sub-circuit
structure (a5.h5 $\to$ a9.h9) only emerges for IOI; SVA and GT are
dominated by infrastructure pathways (layer-0 heads $\to$ MLP0).]

% --- Q/K/V separation ---

\subsubsection{Q/K/V Separation Is Task-Diagnostic}

% REVIEWER: "Is this IOI-specific or general?"
% Answer: It's task-diagnostic. IOI separates, SVA doesn't. The pattern
% tells you about the circuit type, not about the method.

[1 paragraph: Reshaping the 3D tensor into 4D by separating reader
projections (writers $\times$ reader\_heads $\times$ projection $\times$
factors) reveals whether Q, K, and V carry distinct factor subpopulations.
On IOI, the three projection modes are $>$99.9\% pure: each mode loads
on exactly one of Q, K, or V, with a small but consistent Q$\leftrightarrow$K
mixing coefficient (0.020). On SVA, Q and K \emph{do not separate}:
the projection modes mix Q and K at 0.85/0.53. This is structurally
informative: IOI uses distinct factor subpopulations for attention-pattern
(Q/K) vs value-routing (V) because the circuit has separate roles for
``where to look'' vs ``what to match.'' SVA's Q and K are coupled because
subject-verb agreement works through the joint attention score $QK^\top$
--- the same factors help Q find number-bearing subjects and K signal
number information. The Q/K mixing pattern is diagnostic of circuit type:
identity matrix = distinct-role circuit; Q/K merged = joint-attention
circuit. (Appendix~\ref{app:tucker_dimensionality} compares 3D, 4D, and
6D decompositions.)]

% --- Tucker scoring limitation ---

\subsubsection{Scoring vs.\ Interpretability}

[1 paragraph: PCA and Tucker serve complementary roles that happen to
converge on the same decomposition. As a \emph{scoring method}, PCA at
auto-$k$ is universal and requires no tuning; Tucker at fixed low ranks
can beat PCA on specific tasks (IOI: 0.013 vs 0.021) but requires
per-task configuration and fails catastrophically with adaptive ranks.
As an \emph{interpretability tool}, both methods decompose circuits into
the same sub-circuits (cosine 1.000) --- PCA is simpler with one
hyperparameter $k$ vs three ranks. The principled recommendation:
use PCA for both scoring and interpretability; cite Tucker equivalence
as validation that the discovered sub-circuits are method-independent.]


% ============================================================
% 6. EXPERIMENTAL VALIDATION (~2 pages)
% ============================================================

\section{Experimental Validation}
\label{sec:experiments}

% Grokking-style "lines of evidence" opener

[1 paragraph: We validate the factorization and per-factor attribution
with four lines of evidence: (1) null controls (\S6.1),
(2) metric robustness (\S6.2),
(3) circuit faithfulness comparison (\S6.3),
and (4) split-half stability (\S6.4).
All experiments use GPT-2 Small with $N_F = 8192$ (DST
selector, checkpoint atomic-sweep-40).]

% --- 6.1 ---
\subsection{Null Controls}

[1 paragraph: Three null baselines. (i) Random 8D subspaces in
$\mathbb{R}^{8192}$ (20 trials): project scores$_W$, evaluate CMD.
Should be $\gg$0.052 (worse than using all factors). (ii) Shuffled
factors (20 trials): permute factor columns of scores$_W$, run PCA $k=8$,
evaluate CMD. Should be $\approx$0.052 (PCA on destroyed structure
gives no benefit). (iii) Full-factor baseline: CMD = 0.052. PCA's
0.023 is significant if it beats all controls by a wide margin.]

% Results pending from Modal job

% --- 6.2 ---
\subsection{Metric Robustness}

% [FIGURE 6: Grassmannian angles --- NEED]
% Heatmap: 5x5 matrix of mean principal angles between PCA subspaces.
% logit_diff-prob_diff cluster (27 deg), entropy-logit_diff far (82 deg).
% Color scale: blue (0 deg, aligned) to red (90 deg, orthogonal).
\textbf{[FIGURE 6 --- NEED: Grassmannian angle heatmap, 5 metrics.
logit\_diff vs prob\_diff = 27\textdegree\ (aligned),
entropy vs logit\_diff = 82\textdegree\ (orthogonal).]}

[1 paragraph: PCA computed independently with five EAP-IG metrics
(logit\_diff, prob\_diff, gated\_logit\_diff, log\_prob, entropy).
\textbf{Finding}: the IO-specific metrics (logit\_diff, prob\_diff) share
a common subspace (mean angle 27\textdegree) and both yield strong CMD
improvement (2.3$\times$, 2.1$\times$). Global metrics (entropy) produce
nearly orthogonal subspaces (82\textdegree) and no CMD improvement.
This is the correct behavior: the PCA subspace reflects \emph{what the
metric optimizes}, not an artifact of the factorization. The subspace
is stable across metrics that measure the same causal distinction.]

% [TABLE 3: Per-metric CMD --- HAVE]
\textbf{[TABLE 3: Per-metric results. 5 rows (one per metric):
variance explained at $k=8$, full CMD, PCA $k=8$ CMD, improvement ratio.]}

% --- 6.3 ---
\subsection{Circuit Faithfulness}

% [FIGURE 7: Faithfulness curves --- NEED]
% Two panels (or one overlaid):
% Left: faithfulness f(k) vs proportion of circuit kept.
%   Lines: standard EAP, factorized EAP (full), PCA k=8, DAS k=4.
% Right (or inset): CMD/CPR bar chart for each method.
\textbf{[FIGURE 7 --- NEED: Faithfulness curves. Standard EAP vs factorized
EAP (full) vs PCA $k=8$ vs Tucker (10,10,8). x-axis: proportion of circuit
kept. y-axis: normalized faithfulness.]}

[1 paragraph: Standard MIB faithfulness evaluation. Compare methods
at all circuit sizes $k \in [0.001, 1]$: (i) standard EAP (edges ranked
by $|\Delta L_{uv}|$), (ii) factorized EAP with all factors
(sub-edges ranked by $|\Delta L(u,i,v,j)|$), (iii) PCA $k=8$ projected
factorized EAP, (iv) Tucker (10,10,8). Report CMD and CPR
for each. The key comparison is (ii) vs (iii): PCA projection improves
CMD by 2.3$\times$, confirming the denoising effect.]

% --- 6.4 ---
\subsection{Split-Half Stability}

[1 paragraph: PCA computed independently on two disjoint halves of
the IOI prompts (100 each). Principal angles between the half-subspaces
measure stability. Report mean angle at $k=4$, $k=8$, $k=16$.
Expected: $<$15\textdegree\ for top directions confirms the subspace
is a property of the circuit, not prompt-specific overfitting.
Also compare each half to the full-prompt PCA.]

% Results pending from Modal job


% ============================================================
% 7. ATTRIBUTION VS CAUSATION (NEW v6, ~3 pages)
% ============================================================

\section{Attribution and Causation Are Orthogonal in Factor Space}
\label{sec:orthogonality}

% This section is structured to be extractable into a standalone paper
% if desired. It depends on Sections 3-5 (factorization, per-factor
% attribution, sub-circuit decomposition) but the findings are self-contained.

% --- 7.1 Threeway Subspace Comparison ---
\subsection{The Threeway Subspace Comparison}

% SOURCE: SUB_EDGE_DECOMPOSITION_V8.md; outline_ELLIOT_v2.tex lines 183-226
% DATA: Modal fc-results volume /threeway_subspace_comparison/

[Bolded opener: \textbf{The factor subspace that explains the most
attribution variance is nearly orthogonal to the subspace where causal
interventions succeed --- $\sim$88\textdegree\ across 5 tasks.}]

[1 paragraph: PCA, Tucker, and DAS each identify a subspace of the
8192-d factor bank. The first two maximize explained variance in the
attribution tensor; the third maximizes causal interchange accuracy.
We compare all three pairwise via principal angles on the Grassmannian:]

\begin{center}
\small
\begin{tabular}{lcc}
\toprule
\textbf{Pair} & \textbf{Mean cosine} & \textbf{Mean angle} \\
\midrule
PCA vs Tucker & 1.000 & 0\textdegree \\
PCA vs DAS    & $\sim$0.03 & $\sim$88\textdegree \\
Tucker vs DAS & $\sim$0.03 & $\sim$88\textdegree \\
\bottomrule
\end{tabular}
\end{center}

[1 paragraph: PCA and Tucker agree perfectly (as established in
\S\ref{sec:tucker_equivalence}) but both are $\sim$88\textdegree\ from
DAS across IOI, SVA, GT, Gender, Capital. The factors explaining the
most variance in ``how information flows'' are not the factors that
matter for ``what the model computes.'' This is not a failure of PCA
or Tucker --- it is a geometric fact about the relationship between
attribution structure and causal structure in factor space.]

% --- 7.2 Activation-Causal Gap ---
\subsection{The Activation-Causal Gap}

% SOURCE: SUB_EDGE_DECOMPOSITION_V8.md; outline_ELLIOT_v2.tex lines 236-256
% DATA: Modal fc-results volume /activation_causal_gap/

[Bolded opener: \textbf{Factors activate at $r = 0.97$--$0.99$
correlation across tasks but have causal effects at $r = 0.03$--$0.43$,
a 14--40$\times$ gap.}]

[1 paragraph: For each factor, compute activation profile (how strongly
it fires across tasks) and causal profile (how much causal effect per
task). Activation profiles are nearly identical across tasks
($r \geq 0.97$): every factor fires the same regardless of task.
Causal profiles are highly task-specific ($r$ as low as 0.03): which
factors \emph{matter} varies dramatically. This gap is invisible to
scalar edge scores and explains why the attribution subspace (dominated
by high-activation factors) is orthogonal to the causal subspace
(dominated by task-specific factors).]

% --- 7.3 Background/Modulation Decomposition ---
\subsection{Background/Modulation Decomposition}

% SOURCE: SUB_EDGE_DECOMPOSITION_V10.md experiment 2
% DATA: Modal fc-results volume /causal_tucker_modes/

[Bolded opener: \textbf{Tucker C modes capture the modulation
component --- necessary for task performance but not sufficient in
isolation.}]

[1 paragraph: Activation projection experiments using
\texttt{factor\_acts\_enabled()} hooks. Project factor activations
$h = x \cdot F^\top$ onto/orthogonal-to Tucker C subspace at runtime:]

\begin{center}
\small
\begin{tabular}{lcc}
\toprule
$k$ modes & Sufficiency (keep C only) & Necessity (remove C) \\
\midrule
1 & $-$30.4\% (\textbf{reversed}) & 99.6\% \\
2 & $-$29.3\% (\textbf{reversed}) & 97.7\% \\
4 & $-$15.8\% (\textbf{reversed}) & 94.3\% \\
8 & $-$16.8\% (\textbf{reversed}) & 78.9\% \\
\bottomrule
\end{tabular}
\end{center}

[1 paragraph: Removing 8 Tucker C modes drops logit diff by 21\%
(from 0.1956 to 0.1544) --- \textbf{necessary}. But projecting to
\emph{only} Tucker modes reverses the logit diff (negative values =
model predicts the wrong name) --- \textbf{not sufficient}. Like
keeping only the AC component of a signal and discarding the DC offset.
This resolves the activation-causal gap: factors activate universally
because the \textbf{background} dominates activation patterns (baseline
computation sustaining coherent predictions). Causal effects are
concentrated because only the \textbf{modulation} component carries
task-specific information. The 14--40$\times$ gap is the ratio of
background to modulation magnitude.]

% REVIEWER: "This is just saying attribution != causation. That's known."
% Answer: Yes, but we quantify it precisely in factor space and show it's
% not a small effect (~88 deg, not ~45 deg). And we resolve it with the
% background/modulation decomposition, which is new.

[1 paragraph: Note on the factor bank projection null: projecting the
factor bank $F$ (not activations) onto the Tucker C subspace produces
zero effect at all ranks (k=1,2,4,8). This is expected, not a failure:
$F$ is (8192, 768), so any rank-8 subspace in factor-index space maps
to a full-rank-768 output space. The intervention is too weak at the
weight level. The correct intervention is activation projection (above),
which constrains how factors are \emph{read} during forward computation.
Details in Appendix~\ref{app:nullcontrols_extended}.]

% --- 7.4 Per-Counterfactual Stability ---
\subsection{Per-Counterfactual Tucker Stability}

% SOURCE: SUB_EDGE_DECOMPOSITION_V10.md experiment A
% DATA: das_eap/atomic-sweep-40/per_counterfactual_tucker/

[Bolded opener: \textbf{Tucker modes are partially stable across IOI
counterfactual types --- not identical, not orthogonal --- indicating
the structure captures a mix of model-intrinsic geometry and
corruption-specific routing.}]

[1 paragraph: EAP-IG computed separately per IOI counterfactual type
(s2\_io\_flip, s1\_io\_flip, abc, compound). Tucker (10,10,8) decomposed
independently for each. Variance explained varies: s1\_io\_flip is much
more compressible (83.3\%) than compound (39.5\%). Cross-counterfactual
C mode comparison:]

\begin{center}
\small
\begin{tabular}{lccc}
\toprule
\textbf{Pair} & \textbf{Matched cos} & \textbf{Mean angle} & \textbf{Factor Jaccard} \\
\midrule
s2\_io vs compound & 0.440 & 58.0\textdegree & 0.123 \\
s2\_io vs abc & 0.415 & 60.6\textdegree & 0.139 \\
abc vs compound & 0.242 & 73.7\textdegree & 0.072 \\
s1\_io vs abc & 0.240 & 74.1\textdegree & 0.088 \\
s1\_io vs compound & 0.191 & 78.0\textdegree & 0.046 \\
s2\_io vs s1\_io & 0.145 & 78.4\textdegree & 0.111 \\
\bottomrule
\end{tabular}
\end{center}

[1 paragraph: s2\_io\_flip and compound share the most structure
(5/8 modes with cosine $>$0.64, then 3 nearly orthogonal) --- expected
since compound includes s2\_io\_flip. s1\_io\_flip is most different
from all others (angles $\sim$78\textdegree), consistent with its
simpler attribution structure. Factor overlap is low across all pairs
(Jaccard 0.05--0.14) --- different corruptions route through different
factors even when mode directions partially agree. Full tables in
Appendix~\ref{app:per_counterfactual}.]

% --- 7.5 Discrimination ---

\subsection{Factor-Level Discrimination}

% SOURCE: fig2_discrimination.png (from artifacts/spectral_results/neurips_figures/)
% DATA: copied to BATCH2/fig2_discrimination.png

\textbf{[FIGURE --- HAVE: Factor-level discrimination analysis showing
separation between circuit and non-circuit factor usage. Confirms that
circuit-relevant factors form a distinguishable population.
See fig2\_discrimination.png in BATCH2/.]}

[1 paragraph: Factor-level discrimination analysis confirms that
circuit-relevant factors are statistically distinguishable from
background factors. The separation is consistent with the factor overlap
results (\S\ref{sec:factor_overlap}): circuit edges use sparse, disjoint
factor subsets, creating clear signal above the background distribution.]


% --- 7.6 Cross-Task Orthogonality ---

\subsection{Cross-Task Factor Orthogonality}
\label{sec:cross_task_orthogonality}

% SOURCE: cross_task_combined.png (right panel), fig_mode_specificity.png
% DATA: PCA subspace overlap matrix (5 tasks), Tucker mode specificity matrix
% FIGURES: subspace_overlap.png (candidates/), mode_specificity.png (candidates/)

[Bolded opener: \textbf{Different tasks use nearly orthogonal factor
subspaces: mean cosine of principal angles between top-8 PCA subspaces
is 0.03--0.11 across all 10 task pairs.}]

[1 paragraph: For each of 5 tasks (IOI, GT, SVA, Gender Bias,
Capital-Country), we compute the top-8 PCA subspace of the
(writers $\times$ readers $\times$ factors) attribution tensor. Pairwise
subspace overlap (mean cosine of principal angles) is uniformly low:
GT--SVA and GT--Gender are 0.03, the highest pair (IOI--Capital) is
0.11. This is not an artifact of factor sparsity: random 8D subspaces
of $\mathbb{R}^{8192}$ would yield cosine $\sim$0.001. The observed
0.03--0.11 is above random but far below alignment --- each task's
circuit occupies its own region of factor space.]

% [FIGURE: Cross-task subspace overlap heatmap]
\textbf{[FIGURE --- HAVE: 5$\times$5 subspace overlap heatmap (IOI, GT,
SVA, Gender, Capital). Diagonal = 1.0, off-diagonal 0.03--0.11.
RdBu colormap. See subspace\_overlap.png in candidates/.]}

[1 paragraph: Tucker mode specificity confirms this at the component level.
For each task, the best Tucker C mode concentrates 1.5--7.7\% of circuit
energy in its own task's circuit but only 0.3--0.8\% in other tasks' circuits.
SVA modes are the most task-specific (7.7\% own vs 0.3--0.4\% others);
IOI modes are the least specific (1.5\% own vs 0.4--0.7\% others),
consistent with IOI's higher circuit complexity requiring more distributed
factor usage. The specificity matrix is strongly diagonal-dominant:
each task's modes explain its own circuit, not others.]

% [FIGURE: Mode specificity heatmap]
\textbf{[FIGURE --- HAVE: 5$\times$5 mode specificity heatmap showing
\% circuit energy in best mode. Diagonal dominant (1.5--7.7\%) vs
off-diagonal (0.3--0.8\%). See mode\_specificity.png in candidates/.]}

% --- 7.7 Within-Task Sub-Variable Separation ---

\subsection{Within-Task Sub-Variable Separation}
\label{sec:subtask_orthogonality}

% SOURCE: grassmann_heatmap_GOOD.png, 12a_jaccard_mds_v2.png,
%         pca_vs_das_overlap.png, variance_decomposition.png
% DATA: per_variable_das_exploration.json, ioi_subtask_decomposition_v2.json

[Bolded opener: \textbf{Orthogonality holds not only across tasks but
within a single task: IOI sub-variables (IO name, Subject, ABC, Random,
Full flip) route through nearly orthogonal factor subspaces.}]

[1 paragraph: Training separate DAS directions per IOI counterfactual
type at Layer 5 ($k=4$) reveals that the learned causal subspaces are
near-orthogonal. Grassmannian distances between all 10 subtask pairs
range from 0.60 (IO name vs Full flip) to 2.12 (IO name vs Subject),
with 8 of 10 pairs exceeding $\pi/2 \approx 1.57$ (orthogonal).
IO name and Full flip are the closest pair (0.60 radians), which is
expected since Full flip includes the IO name swap. All other pairs
are at or beyond orthogonality.]

% [FIGURE: Grassmannian distance heatmap between IOI subtask DAS subspaces]
\textbf{[FIGURE --- HAVE: 5$\times$5 Grassmann distance heatmap.
IO flip--Full flip = 0.60, all others 1.75--2.12.
$\pi/2 \approx 1.57$ = orthogonal; values $> \pi/2$ = beyond orthogonal.
See grassmann\_heatmap\_GOOD.png in candidates/.]}

[1 paragraph: The factor-level picture confirms this. Each subtask's
DAS direction selects $\sim$17--20 private factors (out of its top 20),
with only 1.0 factor shared on average per pair. MDS embedding of
Jaccard distances shows five well-separated clusters: IO name and
Full flip share 9 factors (thick edge), all other pairs share 0--4.
Out of 82 unique factors identified by DAS across all 5 subtasks,
68 are private to a single subtask and only 14 are shared by 2 or more.
The factor bank provides the vocabulary; each sub-variable selects its
own words.]

% [FIGURE: Jaccard MDS of subtask factor sharing]
\textbf{[FIGURE --- HAVE: MDS of Jaccard distance between IOI subtask
factor sets (DAS top-20, Layer 5). Five colored nodes with edge
thickness $\propto$ shared factors. IO name--Full flip = 9 shared;
most pairs = 0--1 shared. Chance expects $<$1 shared.
See 12a\_jaccard\_mds\_v2.png in candidates/.]}

[1 paragraph: Constrained PCA (unsupervised) and DAS (supervised) find
almost entirely disjoint factor sets: of 20 factors selected per subtask,
mean overlap is 1.0 factor. DAS uses 82 unique factors across subtasks
while CPCA uses only 27 --- the supervised method explores a much larger
region of factor space because it can find subtle causal directions
that don't dominate variance. This reinforces the attribution--causation
orthogonality (\S\ref{sec:orthogonality}): the factors that explain the
most variance (CPCA) are not the factors that carry causal signal (DAS),
and this holds per-subtask, not just globally.]

% [FIGURE: DAS vs CPCA factor overlap]
\textbf{[FIGURE --- HAVE: Three-panel figure. Left: DAS factor overlap
matrix (5$\times$5, 82 unique factors). Center: CPCA factor overlap
matrix (5$\times$5, 27 unique factors). Right: private factors per
subtask (DAS $\sim$17 vs CPCA $\sim$2). Mean overlap 1.0/20 between
methods. See pca\_vs\_das\_overlap.png in candidates/.]}

% REVIEWER: "But orthogonality across subtasks is expected if you
% train separate DAS per type. It's just a property of the optimization."
% Answer: No. DAS finds the direction that maximizes IIA within the same
% 8192-d factor space. Nothing constrains different subtasks to find
% orthogonal directions --- they COULD converge on the same subspace
% if the task used the same factors. Orthogonality is a finding about
% the model's structure, not an artifact of the training procedure.
% The Grassmann distance measures geometry of the converged solution,
% not the optimization path.


% ============================================================
% 8. RELATED WORK (~1 page)
% ============================================================

\section{Related Work}

% Lead with how we differ, not a catalog (Nanda principle)

% REVIEWER: "How does this compare to transcoders? They also decompose
% computations into features."
% Must address. Transcoders decompose MLP computations into SAE features.
% Our factors decompose ALL weight matrices (including attention) through
% a shared bank. Different scope, complementary.

\paragraph{Sparse autoencoders and activation-space decomposition.}
[3-4 sentences: SAEs learn interpretable directions in activation space.
Our factor bank is analogous but in weight space: the model computes
through factors, so they are guaranteed to participate in computation.
SAEs decompose what the model outputs; factors decompose how it computes.
Cite Bricken et al.\ 2023, Huben et al.\ 2023, transcoders.
Factorized DAS with group lasso identifies 9 active factors for IOI
--- but individual factors are gauge-dependent; the causally meaningful
object is the subspace they span.]

% REVIEWER: "The 9 vs 140 comparison is apples to oranges. SAE features
% are at one site; your factors are shared across all sites."
% Acknowledge this. The point is not "we're better than SAEs" but
% "factors provide a different lens."

% NEW v6: SAE comparison paragraph with numbers
SAE features and factors span the same subspace but use different bases.
Projecting 24,576 JumpReLU SAE features (Bloom et al., layer 8, GPT-2
Small) onto the 8,192-factor bank recovers 99.8\% of variance on
average, yet individual alignment is low: the maximum cosine similarity
between any SAE feature and its nearest factor is 0.37, with a mean of
0.15 across all features. This holds regardless of which layer's SAE is
used (layers 5, 7, 8, 9 all yield max cosine $\leq$0.34 with the top
circuit factors). The two decompositions are complementary: SAE features
are monosemantic directions in per-layer activation space, while factors
are shared communication channels in weight space. Critically, because
factors share a global index across all layers, 8\% of factors active in
the IOI circuit participate in two or more pathway stages (top-20 per
edge), rising to 17\% at top-50 --- cross-layer tracking that per-layer
SAE dictionaries cannot provide.
% [FIGURE: SAE cross-layer tracking --- sae_cross_layer_tracking.png]
% See Appendix~\ref{app:figures} Figure~\ref{fig:sae_cross_layer}.

\paragraph{Unsupervised subspace decomposition (NDM).}
[3-4 sentences: Huang \& Hahn (2025) learn an orthogonal rotation of
the residual stream via neighbor distance minimization, achieving Gini
0.71 on IOI subspace patching. Their PCA baseline fails (Gini 0.38)
because PCA on $\mathbb{R}^{d_m}$ activations captures variance, not
task structure. Our PCA succeeds (CMD 0.023 vs 0.052 baseline) because
the factor-attribution tensor \emph{already encodes task relevance} ---
PCA denoises rather than discovers. They decompose \emph{where}
information is stored; we decompose \emph{how} attribution flows.
Their future directions (\S6) propose subspace circuits via weight-space
analysis --- our factorized architecture is exactly this.]

\paragraph{Low-rank structure in weights.}
[2-3 sentences: Kaushik et al.\ (2025) show weight matrices across 500+
LoRA adapters converge to a shared low-rank subspace ($\leq$16 principal
directions). We find analogous low-rank structure in attribution space
within a single model. LoRA's success further supports the low-dimensional
weight-space structure we exploit.]

\paragraph{Circuit discovery methods.}
[2 sentences: ACDC, EAP, EAP-IG, MIB operate at full-edge granularity.
Our contribution decomposes each edge into per-factor contributions and
finds the optimal low-dimensional basis for those contributions.]

\paragraph{Weight-space circuit discovery.}
[2-3 sentences: Tower (2026) shows functional roles (name mover,
S-inhibition, etc.) are recoverable from weight features alone, without
forward passes. Our factorization is complementary: weight features
classify \emph{what role} a head plays; per-factor attribution
decomposes \emph{what information} flows through each edge. Together,
they provide role identification + sub-edge decomposition entirely from
weight structure.]

% NEW v6 related work entries:

\paragraph{Causal abstraction and DAS.}
[2-3 sentences: Geiger et al.\ (2023, 2024). We show the factor
manifold provides effective initialization for DAS. Attribution-based
subspaces are $\sim$88\textdegree\ from the causal subspace --- a
geometric finding, not a method improvement. The orthogonality is
consistent with Geiger et al.'s observation that causal structure
requires interventionist methods, not just observational analysis.]

\paragraph{Conceptor steering and subspace methods.}
[2-3 sentences: Conceptor matrices (NeurIPS 2024) project activations
onto learned subspaces for steering. Factor manifold is a weight-space
analogue. Boolean operations (AND/OR/NOT) on factor-identified task
subspaces produce valid causal interventions (97--100\% IIA).
See Appendix~\ref{app:conceptor}.]

\paragraph{Composition scores.}
[2-3 sentences: Merullo et al.\ (2024) ``Talking Heads'': per-pair SVD
of composition. Our shared factor bank provides a unified cross-edge
basis. Factorized composition via $S_u^\top (F F^\top) S_v$ is
algebraically equivalent to standard composition through the shared
factor bank. See Appendix~\ref{app:merullo}.]


% ============================================================
% 8. DISCUSSION (~0.5 pages)
% ============================================================

\section{Discussion}

% --- What this means ---

[1 paragraph: We adopt the view that sub-edge circuit structure is captured
by factor subspaces --- the spans of factors that circuit projections
selectively activate --- rather than by individual weight matrix ranks or
individual factors. The factor bank makes this precise: PCA on the
per-factor attribution tensor decomposes circuits into interpretable
sub-circuits (IOI: early heads $\to$ MLP0, S-inhibition fan-out, name
movers $\to$ logits), and Tucker decomposition independently recovers
identical factor directions (cosine 1.000, Grassmannian geodesic
$\approx$ 0). This convergence of two mathematically different methods
--- PCA maximizes variance of the flattened matrix, Tucker minimizes
reconstruction error of the 3D tensor --- is strong evidence that the
sub-circuit structure is real, not an artifact of the decomposition
method. Individual factors are basis vectors in the overcomplete
dictionary, analogous to SAE encoder directions --- interpretable as a
catalog, but gauge-dependent; only their spans are mechanistically
invariant. The factor bank also provides a geometric inductive bias for
causal variable discovery: it defines a low-dimensional submanifold of
the Grassmannian where spectral methods succeed (constrained PCA at
98.8\% IIA, zero training) and where DAS optimization converges faster
(91.2\% vs 73.5\% on hard examples). The constraint gradient --- more
geometric constraint yields worse IIA (Riemannian DAS 64\% vs vanilla
99.3\%) --- shows the manifold's value is as initialization, not
ongoing constraint.]

% NEW v6: Orthogonality interpretation

[1 paragraph (NEW v6): The $\sim$88\textdegree\ orthogonality between
attribution and causal subspaces reveals a fundamental distinction.
Attribution captures sensitivity: the directions along which the model's
output changes most when the input is perturbed. Causation captures
computation: the directions through which the model routes task-relevant
information. These are different questions with orthogonal answers in
factor space. The background/modulation decomposition makes this concrete:
factors activate universally to sustain baseline computation (background),
but only a small modulation component carries task-specific signal. Tucker
decomposes the modulation component, which is necessary (21\% logit diff
drop when removed) but not sufficient alone (reversed logit diff when
isolated). The 14--40$\times$ activation-causal gap is the ratio of
background to modulation magnitude.]

[1 paragraph (NEW v6): Each circuit edge uses a distinct factor subset
(cross-pathway Jaccard $<$ 0.5\%), and different counterfactual types
route through partially distinct factor subspaces (principal angles
58--78\textdegree). This suggests the factor bank functions as a
high-dimensional routing table: the bank provides the vocabulary, the
selectors choose the words, and different edges ``say different things''
through the same vocabulary. The sparsity of factor usage per edge
($\sim$20 of 8192) enables this multiplexing without interference.]

% --- Limitations ---

% REVIEWER: "This is a LOT of limitations. Is the paper ready?"
% Be honest. These are real. But the contribution (the factorization
% and adapted methods) is solid even with these limitations.

[1 paragraph: Limitations. (1) The factorized model is an approximation
of the pretrained model, so sub-edge structure exists in the factorized
representation, not directly in the original weights. Model equivalence
(\S3.3) mitigates but does not eliminate this concern.
(2) GPT-2 Small only. Cross-architecture validation (Qwen, Pythia) is
planned but not yet run.
(3) Factorization requires training ($\sim$2 GPU-hours for GPT-2 Small);
per-factor attribution is then free.
(4) Individual factors are gauge-dependent (any rotation of the bank is
equally valid). Only factor \emph{subspaces} (PCA/Tucker mode spans) are
gauge-invariant. The factorization provides the right overcomplete
dictionary to search for subspaces; individual factors are basis vectors,
not interpretable units. This parallels how PCA components aren't
meaningful individually but their span is well-defined.
(5) (NEW v6) Hard-IIA results are on $\sim$34 evaluation examples ---
enough to demonstrate the gradient attractor effect but larger-scale
validation is needed.
(6) (NEW v6) The background/modulation decomposition is demonstrated on
IOI only. Extension to other tasks would strengthen the claim.
(7) (NEW v6) Sub-circuit decomposition recovers known structure
(validation) rather than discovering new biological insights --- the
method is the contribution, the IOI findings are validation.]

% REVIEWER: Limitation (4) is the key one to handle well.
% "You claim per-factor attribution but factors are arbitrary?"
% Answer: The per-factor SCORES are gauge-dependent, but the SUBSPACES
% they identify (via Tucker) are gauge-invariant. The factorization
% gives you the right dictionary to search for subspaces; individual
% factors are basis vectors in that dictionary, not interpretable units.
% The paper's ontological claim: the right unit is the factor subspace,
% not the individual factor.

% --- Future work ---

[1 paragraph: Future directions. (i) Cross-architecture validation:
train factorized Qwen-0.5B, Pythia-160M, and test whether Tucker
structure (core tensor sparsity, mode interpretability) transfers.
(ii) Larger models (GPT-2 Medium/Large).
(iii) Integration with weight-space circuit discovery for end-to-end
role identification + sub-edge decomposition.
(iv) Sparser checkpoints: train factorizations with stronger DST/JumpReLU
penalties to achieve $\sim$500--2000 active factors per head, enabling
pathway-specific factor identification and single-factor knockout experiments.
(v) RAVEL disentanglement: test whether factor subspaces separate
entangled properties (e.g., entity vs.\ attribute).
(vi) Applications: circuit-guided fine-tuning (targeting the active
factor subspace via FactorLoRA), dynamic inference (gating inactive
factors based on selector predictions), conceptor-style compositional
steering via factor subspace algebra (AND/OR/NOT).]


\newpage

% ============================================================
% APPENDICES
% ============================================================

\appendix
\part*{Appendix}

% ============================================================
% APPENDIX A: Gauge Freedom
% ============================================================

\section{Gauge Freedom and QK/OV Distillation}
\label{app:gauge}

% This appendix justifies WHY we distill on composite circuits, not individual W_Q, W_K etc.

\subsection{Rotational gauge freedom in attention}

[State the problem: for any orthogonal $R \in O(d_h)$, replacing
$W_Q \to W_Q R$ and $W_K \to W_K R$ leaves the QK circuit invariant:
$(W_Q R)(W_K R)^\top = W_Q R R^\top W_K^\top = W_Q W_K^\top$.
So individual weight matrices $W_Q, W_K$ are not uniquely determined
by the model's behavior --- only the composite QK = $W_Q W_K^\top$
and OV = $W_V W_O$ are gauge-invariant observables.]

[1 paragraph: This means distilling on individual weights
($\|S_Q F - W_Q\|^2$) is ill-posed: there are infinitely many valid
targets related by arbitrary rotations $R$, and the factorization would
need to ``guess'' the right gauge. Distilling on the composite circuits
($\|QK^{\text{fac}} - QK^{\text{tgt}}\|^2$) eliminates this ambiguity
entirely. The factorization is free to choose any internal gauge for its
$S_Q, S_K$ as long as their product $S_Q F (S_K F)^\top$ matches the
target QK.]

\subsection{Bias reparameterization}

[State the problem: if we zero $b_Q, b_K$ for the factorized model
(since factors should capture bias effects), the attention score changes
because the cross-terms $W_Q b_K$, $b_Q W_K^\top$, and $b_Q \cdot b_K$
are gauge-dependent and nonzero in the pretrained model.]

[1 paragraph: Solution: compute these cross-terms from the \emph{pretrained}
weights before training, store them as frozen buffers
$(\texttt{\_qk\_bias\_q}, \texttt{\_qk\_bias\_k}, \texttt{\_qk\_bias\_const})$,
and add them back during attention score computation. This makes the
factorized attention scores match the pretrained model exactly even with
$b_Q = b_K = 0$. Similarly, $b_V$ is folded through $W_O$ into $b_O$.
See \texttt{reparameterize\_layer\_biases\_} in the codebase.]


% ============================================================
% APPENDIX B: EAP Derivation (EXISTS in Ivan's draft)
% ============================================================

\section{EAP Derivation}
\label{app:eap_derivation}

% Ivan's overleaf_draft.tex Section 2.1 has this COMPLETE.
% Copy from there, lightly edited for notation consistency.

\subsection{Standard EAP}

[COPY FROM Ivan's draft \S2.1 (overleaf\_draft.tex lines 196--292).
The full derivation from $r_u^{(n)} = z_u^{(n)} W_u^{o\top}$ and
$\text{act}_v^{(n)} = \text{LN}(\text{resid}_v^{(n)}) W_v^{in}$ to
the EAP formula with shape annotations. Already publication-ready.]

\subsection{Factorized EAP}

[COPY FROM Ivan's draft \S2.2 (overleaf\_draft.tex lines 404--512).
The four-line algebraic expansion substituting $W = SF$ into the EAP
formula, arriving at $\Delta L_{uv} = \sum_{i,j} \Delta L(u,i,v,j)$.
Already publication-ready. Includes the intuitive description:
``$\Delta L(u,i,v,j)$ answers: how much does writing by module $u$ into
factor subspace $i$, as read by module $v$ through factor subspace $j$,
affect the loss?'']

\subsection{EAP vs.\ EAP-IG}

[1--2 paragraphs:

\textbf{Paragraph 1 (decomposition is exact for both):}
We present the factored decomposition (Proposition~1) for first-order
EAP.  The decomposition extends immediately to EAP-IG by linearity: at
each integration step~$s$ the substitution $W = SF$ is the same
algebra, so the per-factor score at step~$s$ sums to the edge score at
step~$s$.  Averaging over steps preserves the equality:
$\Delta L_{uv}^{\mathrm{IG}} = \sum_{i,j} \Delta L^{\mathrm{IG}}(u,i,v,j)$
where each per-factor EAP-IG score is simply the average of per-factor
scores across integration steps.  No new proof is required.

\textbf{Paragraph 2 (closed-form shortcut is EAP-only):}
What \emph{does not} extend to EAP-IG is the closed-form shortcut
(Appendix~\ref{app:closed_form}), which pre-computes the LayerNorm
Jacobian as a single matrix~$M$.  Under EAP-IG the Jacobian varies at
each interpolation point, so $M$ would need to be recomputed per step.
Our implementation handles this correctly by running the factored
decomposition at each IG step independently via
\texttt{attribute\_rotation.get\_scores\_eap\_ig}.  All scores reported
in the main text use \textbf{factorized EAP-IG} (5 integration steps),
which computes exact per-factor scores at each step and averages.]


% ============================================================
% APPENDIX C: Closed-Form Score (EXISTS in Ivan's draft)
% ============================================================

\section{Closed-Form Factorized EAP Score}
\label{app:closed_form}

% Ivan's overleaf_draft.tex Appendix B has this COMPLETE.
% The LayerNorm Jacobian expansion, M matrix, r encoding, T1/T2 terms, einsum form.

\subsection{LayerNorm Jacobian}

[COPY FROM Ivan's draft Appendix A (overleaf\_draft.tex lines 590--751).
Full derivation of $J_{\text{LN}} = \frac{1}{\sigma}(I - \frac{1}{d_m}
\mathbf{1}\mathbf{1}^\top - \frac{1}{d_m}\hat{x}\hat{x}^\top)$.
Already publication-ready with all intermediate steps.]

\subsection{Expanding $F_i J_{\text{LN},v}^{(n)\top} F_j^\top$}

[COPY FROM Ivan's draft Appendix B \S B.1 (lines 757--829).
Introduction of the Gram matrix $M_{ij} = F_i F_j^\top -
\frac{1}{d_m}(F_i \mathbf{1})(\mathbf{1}^\top F_j^\top)$ and the
factor encoding $r_v^{(n)} = F \hat{x}_v^{(n)}$. Arrives at the
scalar form of the link term.]

\subsection{Final closed-form score}

[COPY FROM Ivan's draft Appendix B \S B.2--B.3 (lines 831--974).
The writer scalar $A_{u,i}^{(n)}$, reader scalar $C_{v,j}^{(n)}$,
substitution chain arriving at
$\Delta L(u,i,v,j) = M_{ij} T_1^{(uv)} - \frac{1}{d_m} T_2^{(uv)}$.
Einsum forms for $T_1$ and $T_2$. Cache contents table.]


% ============================================================
% APPENDIX D: PCA Theorem Proofs
% ============================================================

\section{Proofs of Theorems 1--3}
\label{app:proofs}

\subsection{Theorem 1: Eckart-Young-Mirsky (Optimal Linear Compression)}

[State the theorem formally for our setting. Let $S \in \mathbb{R}^{E
\times N_F}$ be the attribution matrix with SVD $S = U \Sigma V^\top$.
For any rank-$k$ matrix $B$, $\|S - U_k \Sigma_k V_k^\top\|_F \leq
\|S - B\|_F$. Apply: at $k=8$, reconstruction error = 40.8\% of total
variance, which is the minimum achievable by any rank-8 projection.]

\subsection{Theorem 2: Davis-Kahan Subspace Stability}

[State the $\sin\Theta$ theorem. If $S = S_{\text{true}} + E$ and
$\delta_k = \sigma_k - \sigma_{k+1}$ is the singular value gap, then
$\sin\angle(\hat{V}_k, V_k) \leq \|E\|_{\text{op}} / \delta_k$.

Estimate noise level: the Marchenko-Pastur threshold for an
$E \times N_F$ random matrix gives $\sigma_{\text{noise}} = 0.0002$.

Compute bounds:
\begin{itemize}
\item $k=4$: $\delta_4 = 0.0023$, bound = $0.0002 / 0.0023 = 0.087$,
  $\arcsin(0.087) \approx 5.0°$.
\item $k=8$: $\delta_8 = 0.0003$, bound = $0.0002 / 0.0003 = 0.667$,
  $\arcsin(0.667) \approx 41.8°$.
\item $k=16$: $\delta_{16} \approx 0$, bound $\to \infty$, angle $\to 90°$.
\end{itemize}]

\subsection{Theorem 3: Active Subspace Error Bound / SURE Rank Selection}

[If using active subspace (Constantine 2015):
State the Poincar\'{e}-type inequality. The MSE from projecting the
metric onto the top-$k$ subspace is bounded by $\sum_{i>k} \lambda_i /
(\lambda_k - \lambda_{k+1})$. Compute: $k=4$ bound = 14.2, $k=8$
bound = 144.1. The 10$\times$ gap reflects the primary/secondary
eigenvalue separation.]

[If SURE results available:
State Stein's lemma for the truncated SVD denoiser. The unbiased risk
estimate $\text{SURE}(k) = \|S - S_k\|_F^2 + 2 \text{df}(k) \hat\sigma^2$.
The divergence term $\text{df}(k)$ for truncated SVD counts effective
degrees of freedom. Plot $\text{SURE}(k)$ for $k = 1 \ldots 64$, show
minimum at $k \approx$ ?.  This provides non-circular rank selection.]


% ============================================================
% APPENDIX E: Training Details
% ============================================================

\section{Training Details}
\label{app:training}

\subsection{Optimization}

[Optimizer: Adam. Learning rate schedule: reactive halving --- if
$\mathcal{L}_{\text{total}} > 2 \times \mathcal{L}_{\text{prev}}$,
skip the step and halve LR (bounded below by $\texttt{min\_lr}$).
No warmup, no cosine. Only $F$ and $S$ parameters are optimized
(deduplicated by \texttt{id()}); all other parameters frozen from
pretrained model.]

\subsection{Loss components}

[Core loss: $\mathcal{L}_{\text{core}} = \sum_i \|QK_i^{\text{fac}} -
QK_i^{\text{tgt}}\|^2 + \|OV_i^{\text{fac}} - OV_i^{\text{tgt}}\|^2$,
summed over layers.

Auxiliary loss: selector-specific sparsity penalty.
\begin{itemize}
\item L1: $\lambda_{\text{aux}} \cdot \text{mean}(|S|)$
\item JumpReLU: $\lambda_{\text{aux}} \cdot L_0(S)$ (straight-through estimator)
\item DST: $\lambda_{\text{aux}} \cdot |\theta|$ (learnable per-channel threshold)
\end{itemize}

Optional: entropy loss (within-head + across-model deficit),
factor contrastive loss (pairwise $\exp(|\cos| \cdot \text{mult}) - 1$
on learned factors).]

\subsection{Selector types}

% [TABLE A1: Selector comparison --- from slide 14 / existing data]
\textbf{[TABLE A1: Selector types. Type / Mechanism / Auxiliary loss formula / Masking.
L1: soft penalty, continuous.
JumpReLU: per-channel threshold, hard (STE).
JumpReLU-tanh: JumpReLU + tanh-scaled magnitude, hard (STE).
DST: per-channel learnable threshold + BinaryStep STE, hard (STE).]}

\subsection{Hyperparameter ranges}

[$N_F \in \{1024, 2048, 4096, 8192, 16384, 32768, 40000\}$.
$\lambda_{\text{aux}}$ swept per selector type.
At $N_F = 32768$: $\sim$100 active factors per channel ($\sim$0.3\%).
Pareto sweep: trained $>$200 checkpoints across bank sizes and selector
types. Frontier extracted by non-dominated sorting on (CE, L0/channel).]

% [FIGURE A1: Pareto by bank size --- HAVE]
\textbf{[FIGURE A1: Pareto frontier by bank size, all checkpoints.]}

\subsection{Model Equivalence Details}
% [NEW IN V4]
[CREH evaluation methodology. MAS IIA computation. CE gap measurement.
Bidirectional DAS interchange protocol. Details for both GPT-2 and Qwen.]


% ============================================================
% APPENDIX F: Extended Results
% ============================================================

\section{Extended Results}

\subsection{Full CMD/CPR tables}
[CMD and CPR at $k = 1, 2, 4, 8, 16, 32, 64, 128, 256, 512, 1024, 8192$
for PCA, plus full baseline and DAS $k=1,2,4,8$.]

\subsection{Per-checkpoint faithfulness curves}
[Faithfulness curves for all checkpoints shown in the Pareto analysis,
separated by selector type. Show that the PCA improvement is not
checkpoint-specific.]

\subsection{Split-half raw data}
[Full principal angle vectors between half-subspaces at $k=4,8,16$.
Comparison of each half to the full-prompt PCA.]

\subsection{Null control distributions}
[Histograms of CMD for (i) 20 random 8D subspaces, (ii) 20 shuffled-factor
PCA runs. Mark PCA $k=8$ CMD and full baseline CMD on the histogram.]

\subsection{SURE risk curve}
[If computed: $\text{SURE}(k)$ for $k=1 \ldots 64$. Mark the minimum.]

\subsection{Variance explained spectrum}
[Scree plot: singular values $\sigma_1 \ldots \sigma_{64}$ and cumulative
variance explained. Mark the primary/secondary gap at $k=4$ and the
secondary/tail transition at $k=8$.]

\subsection{Cross-Architecture Results}
% Deferred --- Qwen/Pythia experiments not yet run.
% When available: faithfulness curves, Tucker decomposition, CE gap.
% Same format as GPT-2 results for direct comparison.
[DEFERRED to future revision. Cross-architecture validation planned
for Qwen-0.5B and Pythia-160M.]


% ============================================================
% APPENDIX G: Alternative Decomposition Methods
% ============================================================

\section{Alternative Decomposition Methods}
\label{app:alternative_methods}

% Full comparison table showing we tried everything reasonable

\textbf{[TABLE A2: Full method comparison.]}
% Method              | k   | CMD    | Notes
% Varimax k=8         | 8   | 0.019  | Best overall. Orthogonal rotation of PCA.
% PCA k=8             | 8   | 0.023  | Provably optimal linear projection.
% PCA k=16            | 16  | 0.025  | Slightly worse --- noise creeping in.
% ICA k=16            | 16  | 0.032  | Best ICA. Statistical independence weaker than variance.
% PLS-DA n=8          | 8   | 0.045  | Supervised (Wang et al. labels). Worse than unsupervised PCA.
% Group Lasso DAS     | --- | 0.047  | Trained. Supervised.
% EAP-IG baseline     | all | 0.052  | Full N_F-factor baseline (no projection).
% Sparse PCA          | --- | failed | All zeros even at alpha=0.001.
% ICA k=2             | 2   | 0.084  | Too few components.
% ICA k=4             | 4   | 0.080  |

[1 paragraph: We evaluated decomposition methods on the factorized
attribution tensor to verify that PCA/Varimax's performance is not an
artifact of method choice. Table~A2 reports CMD for all methods at their
best $k$. PCA and Varimax dominate; ICA (statistical independence) and
PLS-DA (supervised labels) both underperform unsupervised PCA, confirming
that the low-rank linear structure in the attribution tensor is the
primary signal, and that rotating within the PCA subspace (Varimax)
extracts additional benefit from the sparsity of factor selectors.]

\subsection{ICA}

[2-3 sentences: FastICA at $k = 2, 4, 8, 16$. Best result at $k=16$
(CMD 0.032), worse than PCA at $k=8$ (0.023). ICA seeks statistically
independent components, which is a weaker criterion than variance
maximization for this problem --- the attribution signal is low-rank
and approximately Gaussian in the dominant dimensions, so PCA is
better suited.]

\subsection{PLS-DA}

[2-3 sentences: Partial Least Squares Discriminant Analysis using Wang et
al.\ (2023) IOI circuit labels (13{,}795 circuit edges). Best result at
$n=8$ components (CMD 0.045), better than the full baseline but worse
than unsupervised PCA. The supervised labels help concentrate signal but
PLS overcommits to the binary circuit/background separation at the expense
of fine-grained within-circuit ranking. The unsupervised factor structure
is richer than the binary labels.]


% ============================================================
% APPENDIX H: Stability and Enrichment Analysis
% ============================================================

\section{Stability and Enrichment Analysis}
\label{app:stability_enrichment}

\subsection{Bootstrap Stability of Principal Components}

[Bolded opener: \textbf{The leading PCA components are unstable under
bootstrap resampling, while tail components are rock-solid ---
inverting the normal expectation and explaining the U-shaped CMD curve.}]

[1 paragraph: 50 bootstrap PCA runs (resampling edges with replacement).
Principal angles between each bootstrap subspace and the full-data PCA
subspace: PC0 at 89\textdegree\ (essentially random), PC1 at 85\textdegree,
PC2 at 76\textdegree. But PC6 at 2.7\textdegree, PC7 at 0.2\textdegree.
No individual factor achieves stability $>$ 0.5 (the highest is f44 at
0.32). Interpretation: the dominant variance (PC0--2) is driven by
outlier edges --- high-variance positional/embedding structure that
changes with the edge sample. The robust structure lives in the
smaller components (PC6--7), which capture the stable computational
pathway (induction $\to$ S-inhibition $\to$ name mover). This explains
why $k=8$ works: it captures enough stable structure while the
leading-component noise is diluted by the more stable tail components
in the norm-based ranking.]

% [FIGURE A2: Bootstrap stability --- NEED]
\textbf{[FIGURE A2 --- NEED: Principal angle vs.\ component index for
50 bootstrap runs. Error bars showing spread. Clear two-tier structure:
PC0-3 at 50--90\textdegree, PC4-7 at 0--15\textdegree.]}

\subsection{GSEA Enrichment of Principal Components}

[Bolded opener: \textbf{Gene Set Enrichment Analysis on factor loadings
reveals that duplicate token heads cluster in PC0--1 (6$\times$
enrichment, $p = 0.001$) while induction heads cluster in PC3/PC6.}]

[1 paragraph: We treat each PC's factor loadings as a ranked list and
test enrichment for Wang et al.\ head categories (duplicate token,
induction, S-inhibition, name mover, backup name mover) using standard
GSEA permutation tests. Duplicate token heads show strong enrichment in
PC0--1 (top-100 factors: 6$\times$ enrichment, $p = 0.001$). Induction
heads cluster in PC3 and PC6. Name movers do not show significant
enrichment in any single PC --- they likely use distributed factor
patterns that span multiple components. This categorical structure in
the principal components provides interpretability: PCs are not arbitrary
directions but correspond to functional head categories in the IOI circuit.]

\subsection{Convergent Factors Across Methods}

[1 paragraph: Three factors appear consistently across PLS-DA
and stability analysis: f44 (highest bootstrap stability at 0.32),
f1798, and f4461. These are candidates for ``core IOI factors'' ---
directions in the factor bank that are both computationally important
(high PCA loading) and structurally robust (high bootstrap stability).
However, individual factor identity is gauge-dependent (\S8), so these
should be interpreted as markers of a robust subspace, not as individually
interpretable features.]


% ============================================================
% Appendix: Tucker Dimensionality Comparison
% ============================================================

\section{Tucker Dimensionality Comparison}
\label{app:tucker_dimensionality}

[1 paragraph: We compare three reshapings of the attribution tensor for
Tucker decomposition: 3D (writers $\times$ readers $\times$ factors,
standard), 4D (writers $\times$ reader\_heads $\times$ projection $\times$
factors), and 6D (wr\_layer $\times$ wr\_head $\times$ rd\_layer $\times$
rd\_head $\times$ projection $\times$ factors). All use HOSVD with
randomized truncated SVD on checkpoint atomic-sweep-40, IOI task.]

\textbf{[TABLE A1: Variance explained by decomposition.]}
% Decomposition           | Shape                  | Variance explained
% 3D (10, 10, 8)          | 157 × 445 × 8192       | 59.75%
% 4D (10, 10, 3, 8)       | 157 × 144 × 3 × 8192   | 42.73%
% 6D (4, 6, 4, 6, 3, 8)   | 12 × 12 × 12 × 12 × 3 × 8192 | 40.51%

[1 paragraph: 3D captures the most variance (59.75\%) because writer/reader
modes can freely mix layer/head/projection structure. 4D and 6D impose
structural constraints that reduce total variance capture but yield more
interpretable modes. For CMD scoring, 3D is the default (Section~5.2).
The 4D/6D decompositions are diagnostic tools: they reveal whether the
circuit has distinct Q/K/V subpopulations.]

\subsection{Factor Subspace Equivalence}

\textbf{[TABLE A2: Principal angles between factor mode subspaces.]}
% Comparison   | Mean angle | Spectrum
% 3D vs 4D     | 68.3°      | [2.8, 29.2, 76.2, 82.0, 87.8, 89.0, 89.3, 89.9]
% 3D vs 6D     | 68.3°      | [2.8, 29.3, 76.3, 82.2, 87.9, 89.0, 89.3, 89.9]
% 4D vs 6D     | 0.2°       | [0.0, 0.1, 0.1, 0.1, 0.2, 0.2, 0.3, 0.7]

[1 paragraph: 4D and 6D factor modes span the same subspace (mean angle
0.2\textdegree, max 0.7\textdegree): separating writer layers from heads
does not change which factors matter. 3D factor modes differ substantially
(68.3\textdegree) because the flat writer/reader axes force structural
information into the factor modes. The 3D decomposition conflates
structure and content; 4D/6D separate them. For applications that need
the factor subspace specifically (e.g., ablation experiments), the 4D
decomposition is preferable. For circuit scoring, the 3D decomposition is
preferable because it captures more total variance.]

\subsection{Q/K/V Projection Modes by Task}

\textbf{[TABLE A3: Projection mode purity by task.]}
% Task | Mode 0       | Mode 1       | Mode 2       | Purity  | 4D var%
% IOI  | Q(1.000)     | V(1.000)     | K(1.000)     | >99.9%  | 42.73%
% SVA  | V(1.000)     | K(0.85)/Q(0.53) | Q(0.85)/K(0.53) | Mixed | 13.98%
% GT   | [TBD --- run pending]                                   |

[1 paragraph: On IOI, projection modes separate Q, K, and V with
$>$99.9\% purity (Table A3). A small but consistent Q$\leftrightarrow$K
mixing coefficient (0.020) reflects the attention-score duality $QK^\top$.
On SVA, V separates cleanly but Q and K mix at 0.85/0.53 --- the circuit
uses Q and K jointly for number agreement, so the attribution distributes
symmetrically across both projections. The 4D variance explained drops
from 42.73\% (IOI) to 13.98\% (SVA) because SVA's signal is
MLP-dominated: separating reader projections splits 99\% of the signal
away from the 1\% in attention. The projection separation pattern is
\emph{diagnostic of circuit type}: identity matrix = distinct-role
circuit (IOI), Q/K merged = joint-attention circuit (SVA).]

\subsection{ICA and Varimax}

[1 paragraph: We apply ICA (FastICA) and varimax rotation to the 8 factor
modes from each decomposition. Neither improves sparsity (Gini shift
$<$0.02). The HOSVD modes are already near-optimal --- the useful
structure is in \emph{which} 8 of 8192 factor directions to use, not in
how to rotate within those 8.]


% ============================================================
% Appendix: Tucker Scoring Details
% ============================================================

\section{Tucker Scoring: Full Results and Ablations}
\label{app:tucker_scoring}

% SOURCE: lib/analysis_grassmanian_subspace/factorized_das_eap/tucker_deep_dive/TUCKER_ANALYSIS_COMPLETE.md
% DATA:   Modal volume fc-results, das_eap/atomic-sweep-40/tucker_scoring/
%         das_eap/atomic-sweep-40/tucker_autorank/
%         das_eap/atomic-sweep-40/tucker_adaptive/
%         das_eap/atomic-sweep-40/tucker_cpca/

\subsection{Scoring Aggregation Methods}
\label{app:tucker_scoring_agg}

% SOURCE: TUCKER_ANALYSIS_COMPLETE.md Section 2 "Scoring Methods"
% DATA:   das_eap/atomic-sweep-40/tucker_scoring/{task}/tucker_scoring.json

[1 paragraph: Two aggregation methods for converting Tucker-reconstructed
factor scores into edge scores. \textbf{abs\_of\_sum}:
$|\sum_f \tilde{T}[w,r,f]|$ matches standard EAP-IG's aggregation.
\textbf{sum\_of\_abs}: $\sum_f |\tilde{T}[w,r,f]|$ counts total
unsigned factor contribution. The optimal method varies by task.]

\begin{table}[h]
\centering
\caption{Scoring aggregation comparison (CMD, lower is better).}
\label{tab:tucker_agg}
% SOURCE: TUCKER_ANALYSIS_COMPLETE.md Section 2 "Scoring Methods"
% DATA:   das_eap/atomic-sweep-40/tucker_scoring/{ioi,sva,greater_than}/tucker_scoring.json
\begin{tabular}{llcc}
\toprule
\textbf{Task} & \textbf{Config} & \textbf{sum\_of\_abs CMD} & \textbf{abs\_of\_sum CMD} \\
\midrule
IOI & (10,10,8) & \textbf{0.013} & 0.028 \\
GT  & (20,20,16) & 0.034 & \textbf{0.027} \\
SVA & (15,15,8) & 0.028 & 0.022 \\
\bottomrule
\end{tabular}
\end{table}

[1 paragraph: sum\_of\_abs works better for IOI because the circuit has
many factors contributing with mixed signs along each edge ---
absolute-value-then-sum captures total factor involvement. abs\_of\_sum
works better for GT because the dominant factors are aligned in sign
within each edge, so summation preserves signal while cancellation
removes noise. Both methods should be tested; we report the best for
each task in the main text.]


\subsection{Sparse Factor Modes}
\label{app:tucker_sparse}

% SOURCE: TUCKER_ANALYSIS_COMPLETE.md Section 2 "Sparse Factor Modes"
% DATA:   das_eap/atomic-sweep-40/tucker_scoring/{task}/tucker_scoring.json (sparse configs)

[1 paragraph: We test thresholding each factor mode $U_f[:,k]$ to its
top-$m$ entries by absolute value before Tucker reconstruction. This
removes noise factors while preserving task-relevant ones.]

\begin{table}[h]
\centering
\caption{Effect of sparse factor modes on Tucker CMD.}
\label{tab:tucker_sparse}
% SOURCE: TUCKER_ANALYSIS_COMPLETE.md Section 2 "Sparse Factor Modes"
% DATA:   das_eap/atomic-sweep-40/tucker_scoring/{ioi,greater_than}/tucker_scoring.json
\begin{tabular}{lcc}
\toprule
\textbf{Factor sparsity} & \textbf{IOI CMD} & \textbf{GT CMD} \\
\midrule
Dense (all 8192) & \textbf{0.013} & 0.027 \\
Sparse-200 & 0.018 & 0.018 \\
Sparse-50  & 0.018 & \textbf{0.018} \\
Sparse-20  & 0.018 & 0.028 \\
\bottomrule
\end{tabular}
\end{table}

[1 paragraph: For GT, sparse-50 is better than dense (0.018 vs 0.027) ---
the top 50 factors per mode capture the task-relevant signal and removing
the remaining factors eliminates noise. For IOI, dense is marginally
better (0.013 vs 0.018) because IOI's richer factor structure benefits
from retaining all contributions. The optimal sparsity level correlates
with the factor effective rank (Table~\ref{tab:tucker_effrank} in
\S\ref{sec:analysis}): higher effective rank $\to$ denser modes preferred.]


\subsection{HOSVD vs ALS Refinement}
\label{app:tucker_als}

% SOURCE: TUCKER_ANALYSIS_COMPLETE.md Section 5 "Fast HOSVD (No ALS Refinement)"
% DATA:   das_eap/atomic-sweep-40/tucker_autorank/{task}/tucker_autorank.json (HOSVD-only configs)

[1 paragraph: We compare HOSVD initialization alone (randomized truncated
SVD, no iterative refinement) against HOSVD + Alternating Least Squares
refinement (100 iterations, tolerance $10^{-5}$). ALS refinement is
critical for scoring quality.]

\begin{table}[h]
\centering
\caption{HOSVD-only vs HOSVD+ALS refinement (CMD, lower is better).}
\label{tab:tucker_als}
% SOURCE: TUCKER_ANALYSIS_COMPLETE.md Section 5 "Fast HOSVD"
% DATA:   das_eap/atomic-sweep-40/tucker_autorank/{task}/tucker_autorank.json
\begin{tabular}{lccc}
\toprule
\textbf{Task} & \textbf{HOSVD-only CMD} & \textbf{HOSVD+ALS CMD} & \textbf{Improvement} \\
\midrule
IOI & 0.170 & \textbf{0.013} & 13$\times$ \\
GT  & 0.080 & \textbf{0.018} & 4.4$\times$ \\
SVA & 0.069 & \textbf{0.022} & 3.1$\times$ \\
\bottomrule
\end{tabular}
\end{table}

[1 paragraph: HOSVD-only initialization is catastrophically worse: 13$\times$
worse CMD on IOI, 4.4$\times$ on GT, 3.1$\times$ on SVA. The randomized
SVD provides a coarse starting point, but the alternating least squares
steps are essential for accurate reconstruction of the fine-grained
factor structure. We use HOSVD+ALS (tensorly's \texttt{tucker} with
\texttt{init="svd"}, \texttt{n\_iter\_max=100}) throughout.]


\subsection{Negative Results: Approaches That Did Not Improve Scoring}
\label{app:tucker_negative}

% SOURCE: TUCKER_ANALYSIS_COMPLETE.md Section 5 "What Doesn't Work"
% DATA:   das_eap/atomic-sweep-40/tucker_adaptive/ (asymmetric ranks)
%         das_eap/atomic-sweep-40/tucker_cpca/ (CPCA pre-filtering)

\subsubsection{Asymmetric Tucker Ranks}

% SOURCE: TUCKER_ANALYSIS_COMPLETE.md Section 5 "Asymmetric Ranks"
% DATA:   das_eap/atomic-sweep-40/tucker_adaptive/{task}/

[1 paragraph: Setting mode ranks independently (e.g., $r_w=15, r_r=1,
r_f=8$ for SVA's rank-1 reader structure) produces CMD 0.33--0.81.
Even when a mode has effective rank 1, compressing it to rank 1 in
Tucker removes the 20\% of variance in non-dominant modes that contains
essential edge-discrimination signal. The rank-1 degeneration described
in \S5.1.2 means Tucker \emph{cannot improve rankings} via the dominant
mode, but aggressively compressing the non-dominant modes actively
\emph{destroys} signal.]

\subsubsection{CPCA Pre-filtering}

% SOURCE: TUCKER_ANALYSIS_COMPLETE.md Section 5 "CPCA Pre-filtering"
% DATA:   das_eap/atomic-sweep-40/tucker_cpca/{task}/tucker_cpca_partial.json

[1 paragraph: We tested restricting the factor dimension to factors
identified by constrained PCA (top 1--10\% by selector magnitude)
before Tucker decomposition. This did not help scoring:]

\begin{table}[h]
\centering
\caption{CPCA pre-filtering on SVA (CMD, lower is better).}
\label{tab:tucker_cpca}
% SOURCE: TUCKER_ANALYSIS_COMPLETE.md Section 4 "CPCA Pre-filtering"
% DATA:   das_eap/atomic-sweep-40/tucker_cpca/sva/tucker_cpca_partial.json
\begin{tabular}{lcc}
\toprule
\textbf{Method} & \textbf{Factors retained} & \textbf{CMD} \\
\midrule
Standard EAP-IG (all factors) & 8192 & \textbf{0.015} \\
CPCA filtered 5\% (no Tucker) & 409 & 0.038 \\
CPCA filtered 10\% (no Tucker) & 819 & 0.047 \\
Best CPCA + Tucker & 409--819 & 0.182 \\
\bottomrule
\end{tabular}
\end{table}

[1 paragraph: The top 1\% of factors by selector magnitude holds only
10\% of the EAP-IG attribution energy. The remaining 90\% is distributed
across ``inactive'' factors that still carry real attribution signal.
CPCA pre-filtering works for DAS (\S4.3) because DAS operates at a
single layer where the selector identifies the right causal subspace.
EAP-IG aggregates across all layers and all edges, so the attribution
signal is broadly distributed and pre-filtering removes real signal,
not just noise. This is a fundamental difference between causal
intervention (DAS) and gradient-based attribution (EAP-IG).]

\subsubsection{SVA Edge-Case Methods}

% SOURCE: TUCKER_ANALYSIS_COMPLETE.md Section 5 "Tucker on SVA"
% DATA:   das_eap/atomic-sweep-40/tucker_rank1/ (pending results)

[1 paragraph: No Tucker variant, rank configuration, scoring method,
or pre-processing step improved on SVA's standard EAP-IG CMD of 0.015.
We tested: factor-only PCA ($k=2$ to 512), soft thresholding
($0.5\sigma$ to $5\sigma$), top-$k$ factors per edge ($k=10$ to 1000),
and high-reader-rank Tucker ($r_r=50, 100, 445$). SVA's circuit is
structurally incompatible with Tucker denoising --- it is already clean,
and the rank-1 reader structure means no Tucker configuration can change
edge rankings (see \S5.1.2). The correct interpretation is that SVA's
0.015 CMD is the floor for this checkpoint and task.]


\subsection{Cross-Task Circuit Structure}
\label{app:tucker_crosstask}

% SOURCE: TUCKER_ANALYSIS_COMPLETE.md Section 3 "Cross-Task"
% DATA:   das_eap/atomic-sweep-40/tucker_scoring/{task}/tucker_scoring.json (circuit sub-tensor configs)

\subsubsection{IOI Sub-Circuit Recovery}

% SOURCE: TUCKER_ANALYSIS_COMPLETE.md Section 3 "IOI: Full Circuit Recovery"

[1 paragraph: Tucker (10,10,8) on IOI's circuit sub-tensor (top 1\% of
edges by standard EAP-IG) recovers the complete Wang et al.\ (2022)
IOI circuit automatically. Each head loads on a single writer mode
with loading $>0.99$. The modes are perfectly interpretable:]

\begin{table}[h]
\centering
\caption{IOI Tucker sub-circuits (top 8 core tensor entries, accounting
for 67\% of total energy).}
\label{tab:ioi_subcircuits}
% SOURCE: TUCKER_ANALYSIS_COMPLETE.md Section 3 "IOI: Full Circuit Recovery"
% DATA:   das_eap/atomic-sweep-40/tucker_scoring/ioi/tucker_scoring.json (circuit sub-tensor)
\begin{tabular}{clccl}
\toprule
\textbf{\#} & \textbf{Sub-circuit} & \textbf{Energy} & \textbf{Tier} & \textbf{Known IOI Role} \\
\midrule
1 & a0.h1 $\to$ m0 via F0 & 13.8\% & Primary & Duplicate token $\to$ MLP0 \\
2 & a0.h4 $\to$ m0 via F1 & 11.5\% & Primary & Early attn $\to$ MLP0 \\
3 & a0.h5 $\to$ m0 via F2 & 10.6\% & Primary & Early attn $\to$ MLP0 \\
4 & a0.h3 $\to$ m0 via F3 & 9.7\%  & Primary & Self-attn $\to$ MLP0 \\
5 & a5.h5 $\to$ m5 via F4 & 7.6\%  & Secondary & Induction $\to$ MLP5 \\
6 & a5.h5 $\to$ a8.h6$\langle$v$\rangle$ via F5 & 5.5\% & Secondary & Induction $\to$ S-inhibition \\
7 & a5.h5 $\to$ m5 via F6 & 4.5\%  & Secondary & Induction $\to$ MLP5 \\
8 & a9.h9 $\to$ logits via F7 & 3.8\% & Secondary & Name mover $\to$ output \\
\bottomrule
\end{tabular}
\end{table}

[1 paragraph: Two-tier structure emerges: the \emph{primary tier}
(\#1--4, 46\% of energy) is infrastructure --- layer-0 heads routing
through MLP0, which every task needs for embedding processing. The
\emph{secondary tier} (\#5--8, 21\%) is the actual IOI algorithm:
induction $\to$ S-inhibition $\to$ name mover, the pathway identified
by manual circuit analysis.]

\subsubsection{SVA: Shallow Rank-1 Circuit}

% SOURCE: TUCKER_ANALYSIS_COMPLETE.md Section 3 "SVA: Shallow Rank-1 Circuit"

[1 paragraph: All 8 SVA sub-circuits route through a single reader (m0):]

\begin{table}[h]
\centering
\caption{SVA Tucker sub-circuits (top 5 core tensor entries).}
\label{tab:sva_subcircuits}
% SOURCE: TUCKER_ANALYSIS_COMPLETE.md Section 3 "SVA: Shallow Rank-1 Circuit"
% DATA:   das_eap/atomic-sweep-40/tucker_scoring/sva/tucker_scoring.json
\begin{tabular}{clc}
\toprule
\textbf{\#} & \textbf{Sub-circuit} & \textbf{Energy} \\
\midrule
1 & a0.h5 $\to$ m0 via F0 & 29.2\% \\
2 & a0.h3 $\to$ m0 via F1 & 16.3\% \\
3 & a0.h1 $\to$ m0 via F2 & 14.4\% \\
4 & a0.h4 $\to$ m0 via F3 & 12.5\% \\
5--8 & a0.h\{11,10,6\} + input $\to$ m0 & 24.8\% \\
\bottomrule
\end{tabular}
\end{table}

[1 paragraph: No mid-layer or late-layer structure in the top-1\%
circuit. SVA is a ``shallow'' task: subject-verb agreement is processed
primarily through early attention heads routing to MLP layer 0. The
reader rank-1 structure (Table~\ref{tab:tucker_effrank}) is not just a
scoring diagnostic --- it reflects the circuit's actual architecture.
There is no multi-layer relay; information goes directly from layer-0
attention to MLP0.]

\subsubsection{Greater-Than: MLP Cascade}

% SOURCE: TUCKER_ANALYSIS_COMPLETE.md Section 3 "Greater-Than: MLP Cascade"

[1 paragraph: GT shares IOI/SVA's primary structure (early heads $\to$
m0) but has a distinctive secondary structure --- a 3-layer MLP cascade:]

\begin{table}[h]
\centering
\caption{Greater-Than Tucker sub-circuits (top 10 core tensor entries).}
\label{tab:gt_subcircuits}
% SOURCE: TUCKER_ANALYSIS_COMPLETE.md Section 3 "Greater-Than: MLP Cascade"
% DATA:   das_eap/atomic-sweep-40/tucker_scoring/greater_than/tucker_scoring.json
\begin{tabular}{clcl}
\toprule
\textbf{\#} & \textbf{Sub-circuit} & \textbf{Energy} & \textbf{Role} \\
\midrule
1--6 & a0.h\{1,3,5,4,7,10\} $\to$ m0 & 82\% & Primary: early $\to$ MLP0 \\
7 & a0.h1 $\to$ m1 via F6 & 3.8\% & MLP cascade layer 1 \\
8 & a0.h1 $\to$ m2 via F7 & 3.3\% & MLP cascade layer 2 \\
9 & m0 $\to$ m1 via F6 & 0.5\% & MLP cascade: m0 $\to$ m1 \\
10 & m0 $\to$ m2 via F7 & 0.3\% & MLP cascade: m0 $\to$ m2 \\
\bottomrule
\end{tabular}
\end{table}

[1 paragraph: Factor modes F6 and F7 are \emph{shared} between the
a0.h1$\to$m1, a0.h1$\to$m2, m0$\to$m1, and m0$\to$m2 sub-circuits.
This sharing means the same factor subspace carries the ``number
representation'' through consecutive MLP layers for greater-than
comparison. The MLP cascade (m0$\to$m1$\to$m2) is GT's distinctive
secondary structure, absent in SVA (rank-1 readers) and present in
a different form in IOI (attention relay rather than MLP cascade).]

\subsubsection{Two-Tier Universal Structure}

% SOURCE: TUCKER_ANALYSIS_COMPLETE.md Section 3 "Cross-Task: Universal Primary + Task-Specific Secondary"

[1 paragraph: All three tasks share the same \emph{primary structure}:
layer-0 attention heads $\to$ MLP0 through task-distinct factor
subspaces (F0--F5). This is infrastructure --- embedding processing
that every task requires. The primary tier accounts for 46\% (IOI),
97\% (SVA), and 82\% (GT) of circuit energy. Tasks differ in their
\emph{secondary structure}: IOI has a multi-layer attention relay
(layers 5$\to$8$\to$9), GT has an MLP cascade (m0$\to$m1$\to$m2),
and SVA has no secondary structure at all. The secondary tier is where
task-specific computation happens; Tucker separates it from the
universal infrastructure via the core tensor's block structure.]


\subsection{4D Tucker: Q/K/V Separation}
\label{app:tucker_4d}

% SOURCE: TUCKER_ANALYSIS_COMPLETE.md Section 4 "4D Tucker (Separating Q/K/V)"
% DATA:   das_eap/atomic-sweep-40/tucker_scoring/{task}/ (4D configs)

[1 paragraph: Reshaping the reader axis from flat 445 nodes to
(157 head-components $\times$ 4 Q/K/V/MLP types) yields a 4D tensor.
Tucker on this 4D tensor reveals whether a circuit uses Q, K, and V
independently or jointly.]

\begin{table}[h]
\centering
\caption{4D Tucker Q/K/V projection mode purity by task.}
\label{tab:tucker_4d_purity}
% SOURCE: TUCKER_ANALYSIS_COMPLETE.md Section 4 / Appendix Table A3
% DATA:   das_eap/atomic-sweep-40/tucker_scoring/{ioi,sva}/tucker_scoring.json (4D configs)
\begin{tabular}{lcccc}
\toprule
\textbf{Task} & \textbf{Mode 0} & \textbf{Mode 1} & \textbf{Mode 2} & \textbf{QKV purity} \\
\midrule
IOI & Q (1.000) & V (1.000) & K (1.000) & $>$99.9\% \\
SVA & V (1.000) & K(0.85)/Q(0.53) & Q(0.85)/K(0.53) & Mixed \\
GT  & \multicolumn{4}{c}{[TBD --- pod pending]} \\
\bottomrule
\end{tabular}
\end{table}

[1 paragraph: IOI separates Q, K, V with $>$99.9\% purity --- each
projection type becomes an independent Tucker mode. The small Q$\leftrightarrow$K
mixing (0.020) reflects the $QK^\top$ attention score duality. SVA
separates V cleanly but mixes Q and K (0.85/0.53 loadings), indicating
the circuit uses Q and K jointly for number agreement. This pattern is
\emph{diagnostic of circuit type}: identity-like projection modes indicate
distinct-role circuits (IOI), while Q/K merging indicates joint-attention
circuits (SVA). The 4D variance drops from 42.7\% (IOI) to 14.0\% (SVA)
because SVA is MLP-dominated --- separating attention projections isolates
a minor fraction of the total signal.]


% ============================================================
% REVIEWER CRITIQUE SUMMARY (from v4 reframing)
% ============================================================
%
% Below: every major reviewer concern and how to address it.
% These are the hardest questions. Resolve before writing prose.
%
% ---- CONCERN 1: "You're analyzing a different model" ----
% Q: The factorized model is not GPT-2. How do you know the
%    factor structure isn't an artifact of the factorization?
% A: Model equivalence (Section 3.3). CE gap < 0.02, DAS IIA >= 0.97
%    bidirectional, same circuit edges recovered. The factorized model
%    is constrained (only S and F are free), so it can't invent arbitrary
%    structure --- it must preserve GPT-2's computations to achieve low CE.
% STATUS: Addressed in Section 3.3. Need to strengthen with MAS numbers.
%
% ---- CONCERN 2: "Individual factors are arbitrary (gauge freedom)" ----
% Q: You claim per-factor attribution but factors are basis-dependent.
%    Any rotation gives different per-factor scores. How is this useful?
% A: Individual factors are basis vectors; subspaces are the invariant
%    objects. PCA/Tucker on factor scores identifies gauge-invariant
%    subspaces. The factorization provides the dictionary; downstream
%    analysis extracts the invariant structure. Same as how PCA components
%    individually are meaningless but their span is well-defined.
% STATUS: Acknowledged in limitations. Could be stronger in Section 4.
%    Consider adding a paragraph: "Per-factor scores are basis-dependent,
%    but the subspaces they identify are not."
%
% ---- CONCERN 3: "Tucker only works on IOI" ----
% Q: Tucker results on SVA and GT?
% A: Tucker on the full tensor doesn't scale to weaker circuits.
%    Concentration analysis (PR/k) works on all 3 tasks. Tucker is
%    best used on pre-identified circuit edges, not for discovery.
% STATUS: Acknowledged in Section 5.2. Frame Tucker as interpretability
%    tool, not discovery method.
%
% ---- CONCERN 4: "PCA always finds low-rank structure" ----
% Q: A 70k x 8k matrix will have low-rank structure by definition.
%    What's the null?
% A: Shuffled-factor null: PCA on permuted factors gives CMD = 0.052
%    (no improvement). Random subspaces: CMD >> 0.052. The structure
%    is in the factor attributions, not in the matrix dimensions.
% STATUS: Section 6.1. Need to run and report exact null numbers.
%
% ---- CONCERN 5: "Concentration is trivially explained by score magnitude" ----
% Q: Circuit edges have larger scores. Larger scores = less noise = more
%    concentrated. You're just measuring signal-to-noise, not structure.
% A: RESOLVED by cutting concentration from paper. Magnitude control
%    experiment showed: GT fails entirely, IOI partially structural at
%    high k only, SVA structural but weak. Not robust enough to be a
%    headline claim. Tucker core tensor sparsity replaces this.
% STATUS: CUT from paper. No longer blocking.
%
% ---- CONCERN 6: "Why not just use SAEs?" ----
% Q: SAEs already give per-feature attribution via SAE-patching.
%    What's the advantage of factors?
% A: (1) Factors are shared across all projections --- same dictionary for
%    Q, K, V, O, MLP_in, MLP_out. SAEs are per-layer. (2) Factors are
%    in the model's computation, not post-hoc. (3) No reconstruction
%    error. (4) The shared bank enables cross-layer tracking.
%    BUT: be honest that SAEs are more interpretable per-feature
%    (each SAE feature has a clear activation pattern) while factors
%    are basis-dependent.
% STATUS: Discussed in Related Work. Could be stronger.
%
% ---- CONCERN 7: "Why DAS for equivalence? Why not other probes?" ----
% Q: DAS is one specific method. Model equivalence should be demonstrated
%    with multiple independent methods.
% A: We use THREE levels: CE (behavioral), DAS/MAS (representational),
%    EAP circuit recovery (structural). That's three independent methods
%    at three different levels of analysis.
% STATUS: Section 3.3. This is adequate.
%
% ---- CONCERN 8: "Only GPT-2 Small" ----
% Q: One model is not enough. Cross-architecture?
% A: Acknowledge as limitation. Cross-architecture (Qwen, Pythia) is
%    planned but not yet run. Deferred to future work / revision.
% STATUS: Limitation acknowledged. Future work. Qwen cut from this version.
%
% ---- CONCERN 9: "The factorization requires training" ----
% Q: This isn't training-free. You need to train a factorized model
%    before you can do any of this. That's a significant cost.
% A: True. ~2 GPU-hours for GPT-2 Small. But: (1) the factorized model
%    is reusable across all tasks (train once, analyze many tasks),
%    (2) per-factor attribution is free once the model exists,
%    (3) the Pareto sweep gives guidance on which checkpoint to use.
% STATUS: Acknowledged in limitations.
%
% ---- CONCERN 10: "What did you actually LEARN about IOI?" ----
% Q: You have 8 dimensions, a 4+4 split, head characterization.
%    But what's the NEW insight about IOI that nobody knew?
% A: (1) The primary/secondary split shows infrastructure vs computation
%    are separable in factor space. (2) Each edge only uses 2 of 8 global
%    dimensions --- the circuit is locally low-dimensional.
%    (3) Tucker identifies the specific pathways.
%    Whether these are "new" depends on the audience. For the mech-interp
%    community, the METHOD is new; the IOI findings are validation.
% STATUS: This is the weakest part of the paper. The method is the
%    contribution; the IOI findings are validation, not discovery.
%    Be honest about this framing.
%
% ============================================================


% ============================================================
% FIGURE + TABLE INVENTORY (updated for v5.3)
% ============================================================
%
% MAIN TEXT FIGURES (target: 5-7 figures + 2-4 tables)
%
%   Figure 1  [NEED]  Pipeline overview (standard -> factorized -> per-factor)  PAGE 1
%   Figure 2  [NEED]  Pareto frontier: CE vs L0/channel by selector type        Section 3
%   Figure 3  [HAVE]  PCA layer-aggregated sub-circuit heatmap (IOI, 8 PCs)     Section 5
%                     pca_layer_heatmaps.png --- 13x13 writer layer x reader layer
%   Figure 4  [HAVE]  PCA vs Tucker agreement matrix (cosine + geodesic)        Section 5
%                     pca_tucker_geodesic_agreement.png --- 8x8 |cos| + Gr(1,8192) geodesic
%   Figure 5  [NEED]  Grassmannian angle heatmap across 5 metrics               Section 6
%   Figure 6  [NEED]  Faithfulness curves (std EAP vs fac EAP vs PCA)           Section 6
%
%   Table 1a  [NEED]  PCA and Tucker scoring vs EAP-IG across 5 tasks           Section 5
%   Table 1b  [NEED]  Effective rank per mode across tasks                      Section 5
%   Table 2   [NEED]  Per-metric robustness                                     Section 6
%
% SUPPLEMENTARY FIGURES (all HAVE):
%   pca_subcircuit_zoomed.png     --- IOI top 15 edges per PC with names + values
%   pca_subcircuit_heatmaps_v2.png --- Full 157x445 SymLog-scaled heatmaps (8 PCs)
%   pca_layer_heatmap_ioi.png     --- IOI 8-panel (same as Figure 3)
%   pca_layer_heatmap_sva.png     --- SVA 8-panel
%   pca_layer_heatmap_greater_than.png  --- GT 8-panel
%   pca_layer_heatmap_gender_bias.png   --- Gender 8-panel
%   pca_layer_heatmap_capital_country.png --- Capital 8-panel
%   pca_multi_task_layer_heatmaps.png  --- 5 tasks x 4 PCs combined
%   pca_multi_task_subcircuits.png --- Full heatmaps 5 tasks x 4 PCs
%   pca_multi_task_zoomed.png     --- Zoomed top edges 5 tasks x 4 PCs
%   tucker_layer_heatmaps.png     --- Tucker 8-panel layer heatmap (IOI)
%   tucker_subcircuit_zoomed.png  --- Tucker top edges per factor mode
%   tucker_subcircuit_heatmaps.png --- Tucker full 157x445 heatmaps
%   pca_tucker_agreement.png      --- 8x8 cosine matrix (diagonal = 1.000)
%   pca_vs_tucker_comparison.png  --- Side-by-side matched pairs
%
% APPENDIX FIGURES/TABLES:
%   Figure A1 [HAVE]  Pareto frontier by bank size (all checkpoints)
%   Figure A2 [NEED]  Bootstrap stability: principal angle vs component index    Appendix H
%   Table A1  [NEED]  Selector types comparison                                 Appendix E
%   Table A2  [NEED]  Full alternative method comparison (9+ rows)              Appendix G
%
% ============================================================
\section{Weight-Space Communication Analysis}
\label{app:weight_comms}
% ============================================================

% SOURCE: lib/factorized_spectral/compare_shared_vs_svd_basis.py
%         lib/factorized_spectral/compare_subspace_geodesic.py
% DATA:   atomic-sweep-40 checkpoint (DST, 8192 factors, 95.8\% sparse)

\subsection{Communication Matrices Are Full-Rank}
\label{app:comm_full_rank}

[1 paragraph: For each edge $(u \to v)$, the communication matrix
$C_{uv} = W_O^{u\top} W_V^v$ is $(d_\text{head} \times d_\text{head})$.
We compute singular value concentration $\sigma_1 / \sum_i \sigma_i$
across all IOI edges and random baseline edges. \textbf{Result:}
concentration is 3.9--5.5\% for all pathway types (previous-token
$\to$ induction, induction $\to$ S-inhibition, S-inhibition $\to$
name-mover) and 4.6\% for random edges. Communication is distributed
across all 64 singular values with no dominant channel. This is why
the nuclear norm (sum of all $\sigma_i$) is the appropriate summary
statistic, and why rank-1 dominance tests
capture only $\sim$4\% of the communication bandwidth.]

\subsection{Geodesic Distance Between Communication Subspaces}
\label{app:comm_geodesic}

% SOURCE: compare_subspace_geodesic.py
% DATA:   24 IOI edges (3 pathways) + 8 random baseline edges

[1 paragraph: We compare communication matrix subspaces across edges
via geodesic distance on the Grassmannian $\mathrm{Gr}(k, d_\text{head})$.
For rank-$k$ subspaces (top-$k$ left singular vectors), we compute
principal angles and the geodesic $d_G = \sqrt{\sum \theta_i^2}$.
\textbf{Result:} at all ranks ($k = 1, 3, 5, 10$), within-pathway
geodesic distances are 85--95\% of the maximum possible distance,
with across-pathway distances only slightly larger (discrimination
ratio 1.03--1.07$\times$). Communication subspaces are near-orthogonal
between all edge pairs, regardless of functional relationship.]

[1 paragraph: \textbf{Interpretation:} the geometric structure
lives in the \emph{selector subspaces}, not in the communication
matrices. Selector subspaces show strong task-specific structure:
IOI and SVA route through 88--89$^\circ$ orthogonal subspaces at
task-critical layers (Appendix~\ref{app:selector_orthogonality}).
The communication matrix captures structural \emph{capacity} (the
pipe diameter), while the selector subspace captures functional
\emph{routing} (which pipe to use). These are complementary views.]

\subsection{Factorized vs Standard Composition Scores}
\label{app:fac_vs_std_cs}

% SOURCE: compare_shared_vs_svd_basis.py
% DATA:   24 IOI edges, 7 scoring methods

[1 paragraph: The factorized and standard paths produce numerically
identical scores (geodesic between subspaces $< 0.001$), confirming
that the factorization is faithful.]

[1 paragraph: For cross-edge comparison, individual factor indices
show near-zero overlap (Jaccard@10 $= 0.015$ within-pathway) with
8192 factors ($\sim$344 active per channel). However, the
discrimination ratio (within/across similarity) is $12.7\times$
for factor Jaccard vs $2.0\times$ for SVD cosine. Full
contribution-vector cosine shows within-pathway $\sim$0.75 vs
across-pathway $\sim$0.66.]

[1 paragraph: \textbf{Implication:} the shared vocabulary
operates at the \emph{subspace} level (PCA/Tucker), not
at individual factor indices. This motivates the PCA and Tucker
decompositions as the natural analysis layer above raw factors.]

\subsection{Selector Subspace Orthogonality Across Tasks}
\label{app:selector_orthogonality}

% SOURCE: grassmannian_similarity.pdf (v7 slides, slide 4-5)
% DATA:   atomic-sweep-40, IOI vs SVA selectors per layer
% TODO:   Redo across all 6 tasks with proper controls

[TODO: Redo the Grassmannian orthogonality analysis across all 6
tasks (IOI, SVA, greater-than, hypernymy, capital-country,
gender-bias) with proper controls. Preliminary results (v7 slides):
top-50 factor Jaccard 0.69--0.85 (shared factors), mean principal
angle 75--83$^\circ$ ($\approx$ orthogonal), strengthening to
88--89$^\circ$ at task-critical layers with DST 8192. Tasks share
the same factors but route through orthogonal subspaces.]

% ============================================================
% NEW v6 APPENDICES
% ============================================================


% ============================================================
% Appendix: Per-Counterfactual Tucker Analysis (NEW v6)
% ============================================================

\section{Per-Counterfactual Tucker Analysis}
\label{app:per_counterfactual}

% SOURCE: SUB_EDGE_DECOMPOSITION_V10.md experiment A
% DATA: das_eap/atomic-sweep-40/per_counterfactual_tucker/

\subsection{Variance Explained by Counterfactual Type}

[1 paragraph: Tucker (10,10,8) variance explained varies dramatically
across IOI counterfactual types: s2\_io\_flip 59.1\%, s1\_io\_flip
$\sim$83.3\%, abc $\sim$52.8\%, compound $\sim$39.5\%. s1\_io\_flip
is much more compressible --- simpler attribution structure for that
corruption type. Compound corruptions have the most complex attribution
tensor (least compressible) because they include multiple perturbation
types simultaneously.]

\subsection{Cross-Counterfactual Mode Comparison}

% [TABLE: All 6 pairwise comparisons]
\textbf{[TABLE --- All 6 pairwise comparisons between counterfactual
types. Columns: pair, mean matched cosine, mean principal angle, factor
Jaccard. Ranked from most similar (s2\_io vs compound: 0.440 cos,
58.0\textdegree) to most different (s2\_io vs s1\_io: 0.145 cos,
78.4\textdegree). Include interpretation: s2\_io and compound share
5/8 modes (cosine $>$0.64), then 3 nearly orthogonal; s1\_io produces
structurally simpler tensor with largely distinct modes.]}

\subsection{Interpretation}

[1 paragraph: Different corruption types route through partially
distinct factor subspaces. The \emph{shared} modes likely reflect
model-intrinsic circuit geometry (always activated regardless of
corruption type), while the \emph{distinct} modes reflect
corruption-specific information routing. The s1\_io\_flip case is
instructive: swapping the subject (S1) position rather than the indirect
object (S2) produces a fundamentally simpler attribution structure
(83\% vs $\sim$50\% variance explained) with largely distinct modes,
suggesting S1 and S2 information flows through different factor subspaces
even though both contribute to the same task.]


% ============================================================
% Appendix: Distributed Logit Lens on DAS Dimensions (NEW v6)
% ============================================================

\section{Distributed Logit Lens on DAS Dimensions}
\label{app:logitlens}

% SOURCE: analysis4_logit_lens_distributed/RESULTS.md
%         analysis5_logit_lens_distributed_pt2/NECESSITY_SUFFICIENCY_TAXONOMY.md

\subsection{Per-Dimension Token Categories}

[1 paragraph: Each DAS dimension has a coherent lexical signature when
decoded through the unembedding matrix. SVA dim 0 promotes singular
verbs (IIA 0.920), dim 1 promotes plural pronouns (IIA 0.960). IOI
dims promote names and gender categories. The distributed logit lens
reveals what information each causal dimension carries, providing
interpretability for the otherwise opaque DAS rotation matrix.]

\subsection{Necessity/Sufficiency Taxonomy}

[1 paragraph: SVA: 1 privileged direction with 3 redundant copies.
IOI: inflated 32-dim subspace; greedy forward search finds 5 optimal
dimensions (IIA 0.987) vs full 32 (IIA 0.820, worse due to
over-parameterization). Different tasks have qualitatively different
causal subspace geometry: SVA is rank-1 with redundancy, IOI is
rank-$\sim$5 with padding noise.]

\subsection{JS Divergence Clustering}

[1 paragraph: JS divergence between token distributions decoded from
each DAS dimension reveals semantic clustering. IOI has 3--5
semantically distinct direction types within its 32 DAS dims. Gender
Bias has near-duplicate vectors. This clustering suggests the causal
subspace has interpretable internal structure that factor-constrained
initialization helps discover.]


% ============================================================
% Appendix: Subspace Trimming and k-Sweep (NEW v6)
% ============================================================

\section{Subspace Trimming and Dimensionality Analysis}
\label{app:trimming}

% SOURCE: analysis7_subspace_trimming; analysis9_k_sweep

\subsection{IOI Subspace Trimming}

[1 paragraph: Greedy forward search: 5 dimensions achieve IIA 0.987 vs
full 32-dim DAS at 0.820. $k=32$ is 16$\times$ inflated with $-$0.167
IIA penalty from noise dimensions that interfere with the intervention.
The ``optimal'' DAS dimensionality is much lower than the trained
dimensionality, consistent with the rank-$\sim$5 structure identified
by the distributed logit lens.]

\subsection{Cross-Task k-Sweep}

[1 paragraph: Diminishing returns after $k=4$ on several tasks. $k=4$
near-optimal with improved generalization. Larger $k$ overfits to
training distribution.]

\subsection{Teacher vs Factorized IIA}

[1 paragraph: IOI improves on factorized model (0.820 $\to$ 0.913 at
$k=32$). SVA/Gender/Capital worsen slightly. Factorization enhances IOI
but may subtly damage other tasks --- consistent with the factorization
creating slightly different internal representations that happen to be
better-aligned with IOI's causal structure.]


% ============================================================
% Appendix: Factor Mediation / Causal Discovery (NEW v6)
% ============================================================

\section{Factor Mediation Analysis}
\label{app:mediation}

% SOURCE: lib/causal_discovery/FINDINGS.md

[1 paragraph: PC algorithm on factor activations identifies 286 unique
factors mediating known IOI composition edges (CI-significant at
$\alpha = 0.05$). 90\%+ of mediating factors are unique to a single
edge; 2 factors capture 50\% of mediation. Per-edge factor mediation
is highly specific, consistent with the Jaccard overlap results in
\S\ref{sec:factor_overlap}: the same pattern of factor disjointness
observed in attribution scores is replicated by an independent causal
discovery algorithm (constraint-based PC) operating on factor
activations.]


% ============================================================
% Appendix: Conceptor Steering (NEW v6)
% ============================================================

\section{Conceptor Steering via Factor Subspaces}
\label{app:conceptor}

% SOURCE: experiments/batch6_atlas/06_13_2026/sub_edge_directions/CONCEPTOR_STEERING_WRITEUP.md

[1 paragraph: Boolean operations on DAS-identified factor subspaces:
AND (intersection of two task subspaces), OR (union), NOT (complement)
produce valid causal interventions at 97--100\% IIA. This demonstrates
that factor subspace algebra enables compositional steering without
retraining DAS. The AND operation identifies shared causal directions
between tasks; NOT removes one task's influence while preserving
another's. The factor bank provides the shared basis that makes these
operations well-defined.]


% ============================================================
% Appendix: Merullo Composition Comparison (NEW v6)
% ============================================================

\section{Comparison to Talking Heads Composition Scores}
\label{app:merullo}

% SOURCE: lib/factorized_spectral/MERULLO_COMPARISON.md

[1 paragraph: Merullo et al.\ (2024) compute per-pair SVD of composition
scores between head pairs. Our factor bank provides a global basis: the
Gram matrix $F F^\top$ mediates all cross-layer communication.
Factorized composition via $S_u^\top (F F^\top) S_v$ is algebraically
equivalent to standard composition through the shared factor bank. The
shared basis enables cross-edge comparison that per-pair SVD cannot ---
two edges can be compared in the same factor coordinates, whereas
per-pair SVD produces incomparable singular vectors.]


% ============================================================
% Appendix: Post-Hoc Sparsity Recovery (NEW v6, negative result)
% ============================================================

\section{Post-Hoc Sparsity Recovery}
\label{app:posthoc}

% SOURCE: lib/analysis_grassmanian_subspace/channel_sparsity_posthoc/RESULTS.md

[1 paragraph: Seven unsupervised methods tested for recovering selector
sparsity post-hoc (without access to the selector matrix): activation
thresholding, gradient magnitude, mutual information, correlation
analysis, sparse coding, ICA, and spectral methods. All fail ---
selector sparsity is coordinate-aligned in the learned factor basis
and is unrecoverable from activations or gradients alone without
supervision. The factorization training procedure is essential;
post-hoc analysis cannot substitute. This negative result strengthens
the case for distillation-based factorization: the sparse structure
must be \emph{trained into} the model, not discovered after the fact.]


% ============================================================
% Appendix: Extended Null Controls (NEW v6)
% ============================================================

\section{Extended Null Controls}
\label{app:nullcontrols_extended}

\subsection{Random Subspace Controls}

[20 random 8D subspaces in $\mathbb{R}^{8192}$: CMD distribution.
PCA k=8 CMD should beat all 20 trials by wide margin.]

\subsection{Shuffled Factor Controls}

[20 shuffled-factor PCA runs: CMD distribution. Histogram with PCA
k=8 CMD and full baseline CMD marked.]

\subsection{Sub-Variable Orthogonality Null}

[1 paragraph: Tucker modes from DAS sub-variable attribution tensors
tested against 50 random weightings. Random cosine: 0.08 vs reference;
DAS sub-variables: 0.13. Cross-comparison also not significant.
\textbf{Honest null result} --- the sub-variable Tucker modes do not
carry significantly more information than random weightings. This is
important to report: it constrains the interpretation of Tucker modes
as capturing task-specific (not just data-specific) structure.]

\subsection{Factor Bank Projection Null}

[1 paragraph: Projecting the factor bank $F$ onto Tucker C subspace
produces zero effect at all ranks (k=1,2,4,8). Random subspace
projections also produce zero effect. Explanation: $F$ is (8192, 768),
so any rank-8 subspace of factor-index space maps to a full-rank-768
output space via the overcomplete basis. Weight-level projection is too
weak an intervention; activation-level projection (Section~7.3) is the
correct approach.]


% ============================================================
% Appendix: CMD Evaluation Across 6 Tasks (NEW v6)
% ============================================================

\section{CMD Evaluation of Tucker Factor Subspaces}
\label{app:cmd_6tasks}

% SOURCE: SUB_EDGE_DECOMPOSITION_V10.md experiment 3
% DATA: das_eap/atomic-sweep-40/causal_tucker_modes/

\textbf{[TABLE --- Tucker CMD comparison across 6 tasks. Per-task Tucker
decomposition of each task's own attribution tensor.]}

% Task           | Var Expl | Full CMD | Tucker CMD | FP k=8 | Random k=8 | FP/Rand
% IOI            | 59.2%    | 0.0366   | 0.0178     | 0.0398 | 0.0642     | 0.62
% SVA            | 95.7%    | 0.0116   | 0.0770     | 0.0236 | 0.1032     | 0.23
% greater_than   | 77.1%    | 0.0295   | 0.5164     | 0.1339 | 0.0638     | 2.10
% capital_country| 83.4%    | 0.0209   | 0.0339     | 0.0301 | 0.1118     | 0.27
% hypernymy      | 77.2%    | 0.1011   | 0.1937     | 0.1946 | 0.2903     | 0.67
% gender_bias    | 66.5%    | 0.0173   | 0.0384     | 0.0284 | 0.1259     | 0.23

[1 paragraph: Tucker reconstruction beats full EAP-IG only on IOI
(0.0178 vs 0.0366). Factor-projected beats random on 5/6 tasks
(FP/Rand $<$ 1). greater\_than is the sole outlier --- Tucker
reconstruction catastrophically fails (CMD 0.5164) and random beats
factor-projected. This task may have a qualitatively different
attribution geometry. High variance capture $\neq$ good edge ranking:
SVA achieves 95.7\% variance explained but Tucker reconstruction still
hurts CMD. Interpretation: Tucker identifies a meaningful factor
subspace (not random) but compressing attribution to that subspace
generally loses information needed for optimal circuit ranking. IOI's
improvement was likely a denoising effect specific to that task's
attribution structure.]


% ============================================================
% Appendix: NMI Role Recovery Details (NEW v6)
% ============================================================

\section{Unsupervised Role Recovery Details}
\label{app:nmi_recovery}

% SOURCE: IOI_CIRCUIT_FROM_ATTRIBUTION_DECOMPOSITION.md; TUCKER_CIRCUIT_SECTION_FOR_PAPER.md

[1 paragraph: PCA on the flattened attribution tensor produces per-edge
loading vectors. Clustering these (k-means, $k=5$ matching the number of
Wang et al.\ functional roles) and comparing to ground-truth role
assignments gives NMI = 0.600 ($p$ = 0.004 by permutation test, 1000
permutations). Within-pathway edges have 1.46$\times$ higher cosine
similarity of PCA loading vectors than cross-pathway edges ($p$ = 0.004
by bootstrap). ``Early NMI'' (considering only the first-assigned mode
per head) reaches 0.826. The decomposition recovers known functional
structure without any supervision, circuit labels, or causal model
specification.]


% ============================================================
% v6 FIGURE + TABLE INVENTORY UPDATE
% ============================================================
%
% NEW MAIN TEXT FIGURES (v6 additions):
%   Figure X  [HAVE]  Factor overlap heatmap (4x4 IOI pathway Jaccard)         Section 5
%                     fig3_factor_overlap_stages_v2.png (BATCH2)
%   Figure X  [HAVE]  Per-edge factor fingerprints + Q/K/V sharing             Section 5
%                     fig1_per_edge_factor_fingerprints.png (BATCH2)
%                     fig2_qkv_factor_sharing.png (BATCH2)
%   Figure X  [HAVE]  Factor discrimination analysis                           Section 7
%                     fig2_discrimination.png (BATCH2)
%   Figure X  [HAVE]  Role similarity heatmaps                                 Section 5
%                     25_role_similarity_heatmaps_v2.png (BATCH2)
%
% NEW MAIN TEXT TABLES (v6 additions):
%   Table X: Threeway subspace comparison (PCA vs Tucker vs DAS, 5 tasks)      Section 7
%   Table X: Activation projection sufficiency/necessity (k=1,2,4,8)           Section 7
%   Table X: Per-counterfactual Tucker stability (6 pairwise comparisons)      Section 7
%
% NEW APPENDIX CONTENT (v6 additions):
%   Per-counterfactual Tucker analysis (variance, mode comparison tables)
%   Distributed logit lens (per-dim tokens, necessity/sufficiency taxonomy)
%   Subspace trimming (greedy forward search, k-sweep)
%   Factor mediation (286 CI-significant factors, PC algorithm)
%   Conceptor steering (AND/OR/NOT boolean operations)
%   Merullo composition comparison (shared vs per-pair basis)
%   Post-hoc sparsity recovery (7 methods, all fail -- negative result)
%   Extended null controls (random/shuffled/sub-variable/projection nulls)
%   CMD evaluation across 6 tasks (Tucker vs factor-projected vs random)
%   NMI role recovery details (0.600, p=0.004, bootstrap)
%
% EXISTING FIGURE FILES (v6 additions):
%   docs/for_paper/NEW/figures/sub_circuit_graph/GOOD/good_pca/BATCH2/
%     fig1_per_edge_factor_fingerprints.png
%     fig2_qkv_factor_sharing.png
%     fig2_discrimination.png
%     fig3_factor_overlap_stages_v2.png
%     25_role_similarity_heatmaps_v2.png
%   docs/for_paper/NEW/figures/grassmannian_landscape/
%     grassmannian_landscape_v11_dark_optimal.png
%     grassmannian_landscape_v12_rho.png
%
% ============================================================

% FIGURES THAT EXIST BUT NEED REMAKING:
%   - Pareto plots exist in W&B / slides but need publication-quality versions
%   - Faithfulness curves exist from slides 31-32 but need PCA overlay + cleaner format
%   - Grassmannian angle data exists (Experiment 4) but needs heatmap visualization
%
% FIGURES THAT MUST BE CREATED FROM SCRATCH:
%   - Figure 1 (pipeline overview) --- the most important, defines the paper visually
%
% BLOCKING EXPERIMENTS (must run before paper is writable):
%   1. Null control numbers (random subspaces, shuffled factors)
%   2. Split-half stability numbers
%   3. Faithfulness curves for PCA vs standard EAP-IG across all tasks
%   4. DiM baseline for constrained PCA comparison (is CPCA novel or does DiM match?)
%   5. CPCA-init DAS on SVA and Greater-Than (extend beyond IOI)
%   6. Composition on hard examples (do AND/OR/NOT help where individual methods fail?)
%   7. Hypernymy as zero-shot test of auto-k PCA (not in training set)

% ============================================================
% Appendix: FIGURE GALLERY (NEW v6)
%
% Each figure gets its own entry. Figures can be promoted to
% main text or killed from here. This is the staging area.
% ============================================================

\section{Figure Gallery}
\label{app:figures}

% All figure files are in:
%   docs/for_paper/NEW/figures/sub_circuit_graph/GOOD/good_pca/BATCH2/
%   docs/for_paper/NEW/figures/grassmannian_landscape/
%   artifacts/spectral_results/neurips_figures/
%   artifacts/spectral_results/sae_deep_comparison/
%
% Generation scripts (where available) are in the same BATCH2/ directory.

% --- F1: Per-Edge Factor Fingerprints ---
\subsection{Per-Edge Factor Fingerprints}
\label{fig:per_edge_fingerprints}

\begin{figure}[h!]
\centering
%\includegraphics[width=0.9\textwidth]{figures/sub_circuit_graph/GOOD/good_pca/BATCH2/fig1_per_edge_factor_fingerprints.png}
\caption{\textbf{Per-edge factor fingerprints for Ind$\to$SInh edges.}
Each row is one edge; columns are factor indices. Binary heatmap
showing top-20 factors per edge. Edges between the same head pairs
use largely non-overlapping factor subsets --- each edge has a unique
factor ``fingerprint.'' Generated by \texttt{viz\_factor\_sharing.py}.}
\end{figure}

% STATUS: Main text candidate (Section 5, factor overlap).

% --- F2: Q/K/V Factor Sharing ---
\subsection{Q/K/V Factor Sharing}
\label{fig:qkv_sharing}

\begin{figure}[h!]
\centering
%\includegraphics[width=0.85\textwidth]{figures/sub_circuit_graph/GOOD/good_pca/BATCH2/fig2_qkv_factor_sharing.png}
\caption{\textbf{Q/K/V factor sharing for 4 IOI heads.}
Per head: Venn-style overlap of top factors used by Q, K, and V
projections. Q and K share more factors with each other than with V,
consistent with the QK attention-score duality. Generated by
\texttt{viz\_factor\_sharing.py}.}
\end{figure}

% STATUS: Main text candidate (Section 5, alongside fingerprints).

% --- F3: Factor Discrimination ---
\subsection{Factor-Level Discrimination}
\label{fig:discrimination}

\begin{figure}[h!]
\centering
%\includegraphics[width=0.85\textwidth]{figures/sub_circuit_graph/GOOD/good_pca/BATCH2/fig2_discrimination.png}
\caption{\textbf{Factor-level discrimination: circuit vs non-circuit
factors.} 15.4$\times$ discrimination ratio for circuit factors vs
1.5$\times$ for non-circuit. Circuit-relevant factors form a
statistically distinguishable population. Originally from
\texttt{artifacts/spectral\_results/neurips\_figures/}.}
\end{figure}

% STATUS: Main text candidate (Section 7, attribution-causal gap).

% --- F4: Factor Overlap Heatmap ---
\subsection{Factor Overlap Across Circuit Stages}
\label{fig:factor_overlap}

\begin{figure}[h!]
\centering
%\includegraphics[width=0.6\textwidth]{figures/sub_circuit_graph/GOOD/good_pca/BATCH2/fig3_factor_overlap_stages_v2.png}
\caption{\textbf{4$\times$4 Jaccard overlap between IOI circuit stages.}
Diagonal: within-pathway Jaccard (0.018--0.037). Off-diagonal:
cross-pathway Jaccard (0.001--0.004). Values are mean pairwise
Jaccard between different edges within/across the same stage, NOT
self-similarity. YlOrRd colormap. Generated by
\texttt{viz\_factor\_overlap\_v3.py}.}
\end{figure}

% STATUS: Main text (Section 5, factor overlap).

% --- F5: SAE Cross-Layer Tracking ---
\subsection{SAE vs Factor Cross-Layer Tracking}
\label{fig:sae_cross_layer}

\begin{figure}[h!]
\centering
%\includegraphics[width=0.85\textwidth]{figures/sub_circuit_graph/GOOD/good_pca/BATCH2/sae_cross_layer_tracking.png}
\caption{\textbf{Cross-layer factor reuse in IOI circuit.} Left panel:
8\% of factors participate in 2+ pathway stages at top-20, rising to
17\% at top-50. Right panel: example of a factor active across
Dup$\to$Ind and Ind$\to$SInh stages. Per-layer SAE features cannot
provide this cross-layer tracking. Originally from
\texttt{artifacts/spectral\_results/sae\_deep\_comparison/exp3\_cross\_layer\_tracking.png}.}
\end{figure}

% STATUS: Related Work (SAE comparison) or appendix.

% --- F6: Grassmannian Landscape ---
\subsection{Grassmannian Landscape for DAS Training}
\label{fig:grassmannian}

\begin{figure}[h!]
\centering
%\includegraphics[width=0.48\textwidth]{figures/grassmannian_landscape/grassmannian_landscape_v11_dark_optimal.png}
%\hfill
%\includegraphics[width=0.48\textwidth]{figures/grassmannian_landscape/grassmannian_landscape_v12_rho.png}
\caption{\textbf{Grassmannian landscape for DAS training on hard IOI,
$\mathrm{Gr}(1, 768)$.}
\emph{Left:} IIA as function of angle to DAS optimal ($x$) and
manifold proximity ($y$). On-manifold init (blue) follows high-IIA
ridge to global optimum (star); off-manifold init (orange) trapped.
\emph{Right:} $\rho$ trajectory showing gradient attractor (26\%$\to$83\%).}
\end{figure}

% STATUS: Main text (Section 4.3, already in v5.5).

% --- F7: Role Similarity Heatmaps ---
\subsection{Head-to-Head Role Similarity}
\label{fig:role_similarity}

\begin{figure}[h!]
\centering
%\includegraphics[width=0.85\textwidth]{figures/sub_circuit_graph/GOOD/good_pca/BATCH2/25_role_similarity_heatmaps_v2.png}
\caption{\textbf{Head-to-head similarity: factor profiles vs edge
profiles.} Left: cosine similarity of factor usage profiles (which
factors each head uses). Right: cosine similarity of edge profiles
(which readers each head connects to). Block structure along diagonal
corresponds to Wang et al.\ functional roles (labeled on left axis).
Factor profiles show stronger block structure (higher within-role
similarity) than edge profiles. Generated by
\texttt{viz\_25\_role\_similarity\_v2.py}.}
\end{figure}

% STATUS: Main text candidate (Section 5, role recovery / NMI).

% --- F8: PCA Layer-Aggregated Heatmaps ---
\subsection{PCA Sub-Circuit Layer Heatmaps}
\label{fig:pca_layer_heatmaps}

\begin{figure}[h!]
\centering
%\includegraphics[width=0.95\textwidth]{figures/sub_circuit_graph/GOOD/good_pca/BATCH2/GOOD/pca_layer_heatmaps.png}
\caption{\textbf{Layer-aggregated PCA sub-circuit heatmaps (IOI, 8 PCs).}
Each panel: 13$\times$13 (writer layer $\times$ reader layer). PC0--3:
embed/L0 $\to$ L0 (early processing, 47\% variance). PC4--6: L5 $\to$
L5--L8 (S-inhibition fan-out, 11\%). PC7: L9--L11 $\to$ logits (name
movers, 2\%). Clean three-tier separation.}
\end{figure}

% STATUS: Main text (Section 5, already in v5.5 as Figure 3/4).

% --- F9: PCA-Tucker Agreement ---
\subsection{PCA vs Tucker Factor Direction Agreement}
\label{fig:pca_tucker_agreement}

\begin{figure}[h!]
\centering
%\includegraphics[width=0.6\textwidth]{figures/sub_circuit_graph/GOOD/good_pca/BATCH2/GOOD/pca_tucker_agreement.png}
\caption{\textbf{PCA vs Tucker factor direction agreement (IOI).}
8$\times$8 $|\cos\theta|$ matrix = perfect diagonal. All matched pairs
have cosine $\geq$0.998, confirming identical factor directions from
two mathematically different methods (variance maximization vs tensor
reconstruction). Off-diagonal $<$0.04.}
\end{figure}

% STATUS: Main text (Section 5, already in v5.5 as Figure 4).

% --- F10: MDS Subspace Map ---
\subsection{MDS Subspace Map}
\label{fig:mds_subspace}

\begin{figure}[h!]
\centering
%\includegraphics[width=0.7\textwidth]{figures/sub_circuit_graph/GOOD/good_pca/BATCH2/2_mds_subspace_map.png}
\caption{\textbf{MDS projection of factor subspace geometry.} Each
point is a head; distance reflects Grassmannian distance between
factor subspaces. Functional groups (prev-token, dup-token, induction,
S-inhibition, name-mover) cluster. Generated by \texttt{viz\_mds\_v4.py}.}
\end{figure}

% STATUS: Appendix or supplementary. Good overview but dense.

% --- F11: Factor Loading Heatmap ---
\subsection{Factor Loading Heatmap}
\label{fig:factor_loading}

\begin{figure}[h!]
\centering
%\includegraphics[width=0.85\textwidth]{figures/sub_circuit_graph/GOOD/good_pca/BATCH2/9b_factor_loading_heatmap_v2.png}
\caption{\textbf{Factor loading heatmap.} Per-head loading on each of
the top 8 PCA components. Shows which heads load on which PCs,
confirming the functional group assignments. Generated by
\texttt{viz\_factor\_loading\_heatmap\_v2.py}.}
\end{figure}

% STATUS: Appendix (detailed version of the role assignment).

% --- F12: Cumulative Factor Mass ---
\subsection{Cumulative Factor Mass}
\label{fig:factor_mass}

\begin{figure}[h!]
\centering
%\includegraphics[width=0.7\textwidth]{figures/sub_circuit_graph/GOOD/good_pca/BATCH2/23b_cumulative_factor_mass.png}
\caption{\textbf{Cumulative factor mass for IOI circuit heads.}
Each line is one head; $x$-axis is number of factors (sorted by
magnitude); $y$-axis is cumulative attribution mass. Steep rise
confirms sparse factor usage: most heads reach 80\% mass with
$<$50 factors (of 8192). Generated by
\texttt{viz\_head\_factor\_fingerprints.py}.}
\end{figure}

% STATUS: Appendix (supports factor sparsity claim).

% --- F13: Head Jaccard Heatmap ---
\subsection{Head Pairwise Jaccard Heatmap}
\label{fig:head_jaccard}

\begin{figure}[h!]
\centering
%\includegraphics[width=0.7\textwidth]{figures/sub_circuit_graph/GOOD/good_pca/BATCH2/24_head_jaccard_heatmap.png}
\caption{\textbf{Pairwise Jaccard similarity between IOI heads (top-20
factors).} Block-diagonal structure matches Wang et al.\ roles.
Within-role Jaccard $\sim$0.1--0.2; across-role Jaccard $<$0.05.
Different from the stage-level Jaccard (Figure~\ref{fig:factor_overlap})
which pools edges within pathways. Generated by
\texttt{viz\_head\_factor\_fingerprints.py}.}
\end{figure}

% STATUS: Appendix (supports factor disjointness claim).

% --- F14: Cross-Layer Evolution ---
\subsection{Cross-Layer Factor Subspace Evolution}
\label{fig:cross_layer}

\begin{figure}[h!]
\centering
%\includegraphics[width=0.85\textwidth]{figures/sub_circuit_graph/GOOD/good_pca/BATCH2/5_cross_layer_evolution.png}
\caption{\textbf{Factor subspace evolution across layers.} Grassmannian
distance between layer-$l$ and layer-$l+1$ factor subspaces for IOI
circuit heads. Large jumps at L5 and L9 correspond to functional
transitions (early processing $\to$ S-inhibition $\to$ name movers).}
\end{figure}

% STATUS: Appendix (nice but not essential).

% --- F15: Full 144-Edge Clustering ---
\subsection{Full Edge Clustering}
\label{fig:full_clustering}

\begin{figure}[h!]
\centering
%\includegraphics[width=0.95\textwidth]{figures/sub_circuit_graph/GOOD/good_pca/BATCH2/26_full_144_edge_clustering.png}
\caption{\textbf{Hierarchical clustering of all 144 IOI circuit edges
by factor attribution profile.} Dendrogram with functional role
coloring. Edges from the same functional group cluster together
(NMI = 0.600). The method discovers the circuit's role structure
without supervision. Generated by \texttt{viz\_batch3.py}.}
\end{figure}

% STATUS: Main text candidate (Section 5, role recovery).

% --- F16: CMD Method Comparison ---
\subsection{CMD Method Comparison}
\label{fig:cmd_comparison}

\begin{figure}[h!]
\centering
%\includegraphics[width=0.85\textwidth]{figures/sub_circuit_graph/GOOD/good_pca/BATCH2/14_cmd_method_comparison.png}
\caption{\textbf{CMD comparison across decomposition methods.} PCA,
Tucker, Varimax, ICA, PLS-DA at various ranks. PCA and Varimax
dominate; ICA substantially worse. Generated by
\texttt{viz\_cmd\_method\_comparison.py}.}
\end{figure}

% STATUS: Appendix (Alternative Decomposition Methods).

% --- F17: Scree Plot ---
\subsection{Scree Plot and Cumulative Variance}
\label{fig:scree}

\begin{figure}[h!]
\centering
%\includegraphics[width=0.7\textwidth]{figures/sub_circuit_graph/GOOD/good_pca/BATCH2/15_scree_plot.png}
\caption{\textbf{Scree plot for IOI attribution tensor PCA.}
Primary/secondary gap at $k=4$; secondary/tail transition at $k=8$.
Auto-$k$ at 90\% variance selects $k=108$.}
\end{figure}

% STATUS: Appendix (Extended Results, variance spectrum).

% --- F18: Cumulative Variance by Task ---
\subsection{Cumulative Variance by Task}
\label{fig:cumvar_task}

\begin{figure}[h!]
\centering
%\includegraphics[width=0.7\textwidth]{figures/sub_circuit_graph/GOOD/good_pca/BATCH2/16_cumvar_by_task.png}
\caption{\textbf{Cumulative variance explained by PCA rank, per task.}
SVA reaches 90\% at $k=5$ (simple circuit); Gender requires $k=190$
(complex). The auto-$k$ threshold adapts to circuit complexity.}
\end{figure}

% STATUS: Main text candidate (Section 5, cross-task) or appendix.

% --- F19: DAS Cosine Matrix ---
\subsection{DAS Direction Angles}
\label{fig:das_angles}

\begin{figure}[h!]
\centering
%\includegraphics[width=0.6\textwidth]{figures/sub_circuit_graph/GOOD/good_pca/BATCH2/10_das_direction_angles.png}
\caption{\textbf{Pairwise cosine between PCA components and DAS
directions.} Off-diagonal, confirming the $\sim$88\textdegree\
orthogonality between attribution and causal subspaces.}
\end{figure}

% STATUS: Main text candidate (Section 7, threeway comparison).

% --- F20: IIA vs CMD ---
\subsection{IIA vs CMD Scatterplot}
\label{fig:iia_vs_cmd}

\begin{figure}[h!]
\centering
%\includegraphics[width=0.6\textwidth]{figures/sub_circuit_graph/GOOD/good_pca/BATCH2/17_iia_vs_cmd.png}
\caption{\textbf{IIA vs CMD across methods and tasks.} Shows the
tradeoff between causal accuracy (IIA) and circuit faithfulness (CMD).
Points colored by method. PCA-projected factorized EAP achieves the
best CMD; trained DAS achieves the best IIA.}
\end{figure}

% STATUS: Appendix (nice summary but not essential).


\end{document}
